\documentclass[parskip=full]{article}

%%%%%%%%%%%%%%%%%%%%%%%%%%%%%%%%%%%%%%%%%%%%%%%%%
% Basic includes
%%%%%%%%%%%%%%%%%%%%%%%%%%%%%%%%%%%%%%%%%%%%%%%%%
\usepackage{array}
\usepackage{makecell}
\usepackage{amsmath}
\usepackage{amsfonts}
\usepackage{verbatim}
\usepackage{mathrsfs}
\usepackage{amssymb}
\usepackage{amsthm}
\usepackage{bbm}
\usepackage{bm}
%\usepackage[capposition=top]{floatrow}
\usepackage[dvipsnames]{xcolor}
\usepackage{cancel}
\usepackage{ulem}
\usepackage[labelfont=bf]{caption}
\usepackage{subcaption}
\usepackage{graphicx}
\usepackage{float}
\usepackage{natbib}
\usepackage{booktabs}
\usepackage{longtable}
\usepackage{multirow}
\usepackage{rotating}
\usepackage{pdflscape}
\usepackage[title]{appendix}
\usepackage{footmisc}
\usepackage{bigdelim}
\usepackage{tabularx}
\usepackage{xspace}
\newcommand{\A}{\ensuremath{\mathcal{A}}\xspace}
\newcommand{\B}{\ensuremath{\mathcal{B}}\xspace}
\newcommand\pa[1]{\ensuremath{\left(#1\right)}}
\usepackage{epigraph}
\usepackage{setspace}
%\onehalfspacing
\setstretch{1.5}

%%%%%%%%%%%%%%%%%%%%%%%%%%%%%%%%%%%%%%%%%%%%%%%%%%%
% Microtype
%%%%%%%%%%%%%%%%%%%%%%%%%%%%%%%%%%%%%%%%%%%%%%%%%%%
\usepackage[activate={true,nocompatibility},final,tracking=true,kerning=true,spacing=true,factor=1100,stretch=10,shrink=10]{microtype}
% activate={true,nocompatibility} - activate protrusion and expansion
% final - enable microtype; use ``draft'' to disable
%tracking=true, kerning=true, spacing=true - activate these techniques
% factor=1100 - add 10% to the protrusion amount (default is 1000)
% stretch=10, shrink=10 - reduce stretchability/shrinkability (default is 20/20)

%%%%%%%%%%%%%%%%%%%%%%%%%%%%%%%%%%%%%%%%%%%%%%%%%%%
% Hyperref setup
%%%%%%%%%%%%%%%%%%%%%%%%%%%%%%%%%%%%%%%%%%%%%%%%%%%
\usepackage{url}
\makeatletter
\g@addto@macro{\UrlBreaks}{\UrlOrds}
\makeatother
\usepackage{hyperref}
\hypersetup{
  colorlinks,%
  citecolor=black,%
  filecolor=black,%
  linkcolor=blue,%
  urlcolor=black
}

%%%%%%%%%%%%%%%%%%%%%%%%%%%%%%%%%%%%%%%%%%%%%%%%%%%
% Font setup
%%%%%%%%%%%%%%%%%%%%%%%%%%%%%%%%%%%%%%%%%%%%%%%%%%%
\usepackage{mathpazo} % math & rm
%\linespread{1}        % Palatino needs more leading (space between lines)
\usepackage{parskip}
\usepackage[scaled]{helvet} % ss
\usepackage{courier} % tt
\usepackage{natbib}

\usepackage{filecontents}

\normalfont
\usepackage[T1]{fontenc}

\usepackage{titlesec}
\titleformat{\section}
{\normalfont\fontsize{12}{17}\bfseries}
{\thesection}
{1em}
{}

\titleformat{\subsection}
{\normalfont\fontsize{11}{16}\bfseries}
{\thesubsection}
{1em}
{}

%%%%%%%%%%%%%%%%%%%%%%%%%%%%%%%%%%%%%%%%%%%%%%%%%
% macros
%%%%%%%%%%%%%%%%%%%%%%%%%%%%%%%%%%%%%%%%%%%%%%%%%
\newcolumntype{C}[1]{>{\centering\let\newline\\\arraybackslash\hspace{0pt}}m{#1}}
\newcommand{\indicator}[1]{\mathbbm{1}_{\left\{ {#1} \right\} }}
%\newcommand{\resultsloc}{../code/python/output/}
\newcommand{\subfigwidth}{0.475\textwidth}
\newcommand{\subfigwidthtwo}{0.375\textwidth}
\newcommand{\figwidth}{0.6\textwidth}
\DeclareMathOperator*{\argmax}{arg\,max}

\setlength\epigraphwidth{0.9\textwidth}
\setlength\epigraphrule{0pt}
\renewcommand{\epigraphflush}{center}
\usepackage{etoolbox}
\makeatletter
\patchcmd{\epigraph}{\@epitext{#1}}{\itshape\@epitext{#1}}{}{}
\makeatother

%%%%%%%%%%%%%%%%%%%%%%%%%%%%%%%%%%%%%%%%%%%%%%%%%
% Page and paragraph settings
%%%%%%%%%%%%%%%%%%%%%%%%%%%%%%%%%%%%%%%%%%%%%%%%%
\setlength{\parskip}{0.0pc} \setlength{\parindent}{0.25in}
\setlength{\topmargin}{-3pc}
\setlength{\textheight}{9in}
\setlength{\oddsidemargin}{0pc} \setlength{\evensidemargin}{0pc}
\setlength{\textwidth}{6.5in}
\makeatletter
\@fpsep\textheight
\makeatother

\newenvironment{fignote}{\begin{quote}\footnotesize}{\end{quote}}
\def\sym#1{\ifmmode^{#1}\else\(^{#1}\)\fi}
\renewcommand{\baselinestretch}{1.5}

% Formats the notes at the bottom of tables and figures.
\newenvironment{fignote2}[1]
{
\par
\medskip
\begin{minipage}{#1}
\footnotesize
}
{
\end{minipage}
}

% Formats the notes at the bottom of tables and figures.
\newenvironment{fignote3}[1]
{
%\par
\begin{minipage}{#1}
\footnotesize
}
{
\end{minipage}
}

%%%%%%%%%%%%%%%%%%%%%%%%%%%%%%%%%%%%%%%%%%%%%%%%%
% Title page
%%%%%%%%%%%%%%%%%%%%%%%%%%%%%%%%%%%%%%%%%%%%%%%%%
%%%%%%%%%%%%%%%%%%%%%%%%%%%%%%%%%%%%%%%%%%%%%%%%%
% Title page
%%%%%%%%%%%%%%%%%%%%%%%%%%%%%%%%%%%%%%%%%%%%%%%%%

\author{
  Shahar Rotberg\footnote{University of Toronto. Contact: \url{shahar.rotberg@mail.utoronto.ca}.} \and Lin Zhang\footnote{CMHC. Contact: \url{lzhang@cmhc-schl.gc.ca.}. The opinions expressed herein do not represent the views of CMHC. We thank Joseph B. Steinberg for kindly letting us use his high-powered computer for this paper.}
}

\title{\vspace{-1.8cm}A General Equilibrium Analysis of Rent Control with Heterogeneous Households}

\pagestyle{plain}

\begin{document}

\maketitle

\setstretch{1.5}
\begin{abstract}
\onehalfspacing
\noindent Does rent control improve housing affordability and societal welfare? To answer these questions, we model rent control in a general equilibrium model with heterogeneous households sorting into rental markets and choosing tenancy duration. We discipline the model with data from Toronto where rental markets are partly controlled and partly uncontrolled. In the long run, full decontrol would raise average rents, but raise welfare through reducing rents of young, low-income households. Results largely hold in transition. No degree of regulation is optimal in the long run. Full rent control would reduce the quantity of rental housing supplied and raise average rents.
\\[1.5ex]
% words=100

\noindent \textbf{JEL Classification:} E17; R31; R52. 

\noindent \textbf{Keywords:} General equilibrium, rent control, housing markets.
\end{abstract}
\setstretch{1.5}

\section{Introduction}

Despite its unpopularity among economists, rent control has been adopted in many major cities, and has evolved from rent freeze to exemptions to tenancy deregulation \citep{arnott-1995,arnott-2003}. Economists have provided partial equilibrium analyses on rent freeze and rent control with exemptions \citep[among many others]{hayek-1929,fs-1946,fallis-1985}. However, tenancy deregulation (i.e., rent control), which imposes rent controls while tenants remain in the same unit but allows landlords to reset rents when tenants change, has been given scant attention by the general equilibrium literature. 

This paper aims to fill this gap by modelling rent controls with exemptions and tenancy deregulation in a general equilibrium model with heterogeneous households. In the model, rent-controlled and uncontrolled (i.e., free) markets  operate simultaneously. We calibrate the model to data from Toronto where part of the rental market under rent control and part is free, and use the model to explore rent control's impact on housing affordability and societal welfare, and its distributional consequences for different households.

Households in our model make endogenous tenure and tenancy choices, and are sorted into rental markets. Households face mandatory retirement after a certain age, and working-age households supply labor inelastically. Households choose whether to rent or to own. Households who choose to own a house have access to long-term mortgage contracts.\footnote{Mortgage interest payments cannot be deducted from taxable income because this policy does not exist in Canada.} Renters choose whether to rent in the rent-controlled market or the free market. Controlled rents decrease over time in real terms, and thus the rents new tenants pay in the rent-controlled market are higher than the rents they would initially pay in the free market. A consequence of this is that higher-income households sort into the free rental market, as they stay only a short time in their units before buying homes, and thus do not benefit from rent control. This assumption is strongly supported by data from Toronto.\footnote{Toronto data indicate that there is sorting in the rental market; higher-income households tend to sort into the free market. 2016 Toronto Census data show that the average household income of renters who rent in the free market is 20\% higher than the average household income of renters who rent in the rent-controlled market. Moreover, the average rental unit in the decontrolled market in Toronto is 24\% better than the average rental unit in the regulated market \citep{zhang-gou} along dimensions such as location, size, quality, etc. Our model replicates these facts quite closely. Furthermore, sorting is likely quite realistic due to the visible nature of rent control laws in reality, which normally state that a unit is under rent control only if it was built before a certain year. Since the year of construction is readily available information for households in reality, they can know ahead of time whether they will be living in a unit that is under rent control or not.}

Rental housing is supplied by a competitive rental company with convex marginal costs. Similarly to \citet{fallis-1985}, we model the rental housing supply after Toronto's rental market structure, where new construction is not subject to rent control but old construction is rent-controlled. This means that the rental housing supply depends on the marginal rent, which is the free market rent, not the average rent or rent in the rent-controlled market. A competitive housing construction sector with convex marginal costs supplies new housing, which are added to the depreciated housing stock. The government imposes proportional property taxes and progressive taxes on labor and capital incomes. With its proceeds, the government pays pensions to retired households, lump sum transfers to all households, and provides a public good that does not enter household utility. 

We calibrate the model to match moments in the data for Toronto. For instance, as of 2016,  12.4\% of the market in Toronto is decontrolled and the rest is under rent control. Among these moments, one crucial moment for our results is the fraction of renters that spend over 30\% of their income on rent. We estimate this data point using data from the 2016 wave of the Census and use this statistic to calibrate the smallest dwelling available for rent. This calibration in turn affects the degree to which changes in rents affect low-income households' non-housing consumption and their marginal utility. Since low-income households have the highest marginal utilities, changes in their welfare can substantially affect overall societal welfare. In section \ref{subsection:non_targeted_statistics} we show that the model performs well on important aggregate, life-cycle, and housing tenure statistics that we do not target in our calibration.

Using our calibrated model, we explore the implications of rent control for equilibrium rents, house prices, societal welfare, and distributional outcomes. We follow common practice in the literature and define societal welfare as the ex-ante expected utility of newborn households \citep{conesa-et-al,gkkoc}. We use our calibrated model with partial rent control as the benchmark we compare our experiments to. In all of our experiments we keep the tax system constant and restore fiscal balance using lump sum transfers. In our first three experiments we compare long-run stationary equilibria, while in our last experiment we compute the transition dynamics from one long-run equilibrium to another long-run equilibrium. 

Our first experiment highlights the impacts of decontrolling rental markets. We compute the long-run equilibrium in our model under the extreme case of fully decontrolled rental markets and compare it to the benchmark equilibrium. We find that the main tension that rent decontrol brings about is that between average housing affordability and housing affordability early in life. Decontrolling would reduce the value of staying longer in the same rental unit, and thus encourage renters to save more to buy a house, which would reduce total rental demand, and thus reduce free market rents. This would reduce initial rents paid by tenants when moving into a rental unit by 6.7\%, but raise the average rent in the economy by 4.1\%. The lower rents paid by tenants moving into a new unit mostly benefit low-income, low-asset young households. Because these households have the highest marginal utilities, their large welfare gains would also raise societal welfare by 0.26\%. Homeownership would also rise by 0.65p.p. (percentage points) because staying longer as a renter is less beneficial when rental markets are decontrolled. Given the increase in homeownership, house prices would increase by 0.8\% to clear the housing market. As a result, property tax revenue would increase, enabling the government to provide lump sum transfers to households equal to 0.03\% of the average labor income. Lastly, the rise in homeownership would reduce the overall demand for rental housing by 0.7\%, and thus free market rents would decline by 3.3\%. 

Our model shows that the impact of rent decontrol on welfare would be unequally distributed across households. Newborn households with the lowest labor productivity would gain substantially, because early in life, when their income and assets are low, their rents would be lower. This would increase their disposable income for non-housing consumption during years in which they have little resources, and thus raise their welfare. Newborn households with low-to-mid-productivities, on the other hand, put more value on average rents than the rents early in life, and thus would lose. This is because they expect to stay long enough in the same controlled rental unit to benefit from rent control, but have higher incomes early in life such that the decline in rents early in life would affect their non-housing consumption less significantly. Newborn households with high productivities would gain. Because they have higher incomes, they expect to become homeowners quickly, and thus they sort into the free market. Because free market rents would decline by 3.7\% they would benefit from lower rents. 

Lastly, we measure the various sources of welfare gains. Since the impact on equilibrium rents would be the largest and affect low-income households the most, rents would have a first order impact on welfare, but as we discussed above, the welfare impact would be heterogeneous. The welfare effect of the lump sum transfers, in contrast, would be positive across the board, but have a second order effect on welfare. The increase in house prices would have a third order effect on welfare because it would mostly impact high income households whose marginal utilities are lowest. This effect would be negative across the board because it would causes households to delay their homeownership and reduce their house size once they buy.

In our second experiment, we compare the benchmark to the long-run equilibrium in which all rental units are under rent control. Full rent control would encourage renters, who rented in the free market in the benchmark, to stay longer in the same unit. However, since these renters would move less, they would also demand units that are smaller relative to what they demanded in the benchmark. These two effects would reduce overall housing demand, and thus reduce house prices by 0.3\%. With lower house prices, more renters would be able to purchase a house, and homeownership would rise by 0.3p.p. However, the elimination of the free market would also reduce the marginal rent by 6.77\%, because the free market rent in the benchmark is higher than the average rent in the rent-controlled market. Because rental housing supply is determined by the marginal rent, this would reduce the quantity of rental housing supplied by 1.4\%, and thus increase initial rents paid by new tenants by 0.1\% as well as average rents by 0.4\%. The latter two effects would hurt virtually all households and reduce welfare by 0.5\%. The key takeaways from these results are that they support the claim that a fully rent-controlled market reduces the quantity of rental housing supplied in equilibrium and increases average rental prices, and that exemptions mitigate these effects. 

In our third experiment we search for the optimal degree of rent control. We increase and decrease the share of rental housing that is mandated by the government to be under rent control and compute welfare under each case in the long-run equilibrium. Our model shows that no degree of rent control is optimal, implying that full rent decontrol is welfare-maximizing. Why? This is because as the rent-controlled market share declines, sorting of households into rental markets implies that equilibrium rents in the rent-controlled market must increase in order to shift more renters to the free market. Since lower-income renters tend to stay longer in their rental units and benefit more from rent control, they would end up paying higher equilibrium rents the lower the share of the rent-controlled market gets. Since low-income renters have the highest marginal utilities, this negative effect on them would also substantially negatively impact overall welfare. 

In our last experiment, we compute the cost of transitioning from the benchmark equilibrium to the long-run equilibrium with full rent decontrol, or the optimal long-run equilibrium. We model this as an fully unexpected shock, in which the government fully decontrols rental markets permanently. Rents would rapidly decline upon implementation, resulting in welfare rising immediately and remaining positive almost all along the transition path. The main reason for the instantaneous reduction in rent is that renters who live in large rental units and who were previously protected by rent control, would no longer find it optimal to remain in those large units. This would cause overall rental housing demand to decline, rents to decrease, and low-income, high marginal utility newborn renters to benefit substantially even upon implementation, raising also the ex-ante expected welfare of newborns.

Our examination here builds on the efforts by others to model and study rent controls. Some prominent economists such as \citet{hayek-1929} and \citet{fs-1946} study how rent freeze negatively affects rental supply, while economics textbooks analyze its deadweight loss impact. \citet{fallis-1985} use a static microeconomic model of rental markets. Their model includes uncontrolled and controlled markets interacting with each other and predicts that uncontrolled rents are higher than controlled ones. Because their model does not include homeownership, it does not account for the impact of tenure decisions on equilibrium prices. Moreover, because their model is static, it does not account for impact on decisions of the dynamic nature of rent control; the longer a renter stays in their unit, the more they benefit from rent control. As shown in \citet{diamond-2019}, empirical studies have moved ahead of theoretical modelling in studying tenancy deregulation with exemptions or partial rent control and thereby require the equilibrium analysis of rent control with exemptions, especially in order to understand who the winners and losers are. We model partial rent control in a general equilibrium model with heterogeneous households and examine how rent control affects housing affordability and the welfare of different households.

The work in this paper is related to \citet{fmn-2022}, who use a heterogeneous-agent housing general equilibrium model, to study the impact of rent vouchers, rent subsidies, and rent control, on housing affordability. Their rent control set up in the paper, however, rent-controlled units in \citet{fmn-2022} are similar to a sub-segment of social and affordable housing in Canada. This is because in their model renters living in controlled units enjoy a rent discount relative to the free market immediately upon entering their unit. In contrast, in our model renters sort into rental markets, which means that to clear rental markets renters living under rent control initially must pay a premium to rent in the controlled market. Our work is also different as it focuses on the feedback effects between the rent-controlled and free markets in the presence of sorting, and the transition out of rent control.

Our work provides policy implications for transitioning out of rent controls. While most of studies are focused on the misallocation of resources due to rent controls \citep{gl-2003}, to the best of our knowledge there are no theoretical studies on the transition out of rent control. \citet{sims-2007} empirically evaluates the effects of rent control on rental prices and rental supply using the sudden stop of rent control in Massachusetts, USA. Similarly to what we show with our model, he also finds that rent decontrol increases average rents. In contrast to \citet{sims-2007}, however, because we use a model in which agents make optimal decision to optimize their utility, we are able to perform a welfare analysis and a compute what the transition from the benchmark with rent control to a fully decontrolled rental market would look like. %Contrary to common beliefs, our model shows that rent control hurts young, low-income renters rather than help them, and that the cost of transitioning from rent control to rent decontrol, which some fear would be large, would actually be minimal. 

The paper is organized as follows. Section \ref{sec:insitutional_background} provides institutional background of rent control in Toronto. Section \ref{sec:model} lays out the model environment. Section \ref{sec:calibration} presents model calibration the non-targeted statistics. Section \ref{section:experiments} presents models results and transitions from rent control to no control, and section \ref{sec:conclusion} concludes. 

\section{Institutional background on rent control in Toronto, Canada}\label{sec:insitutional_background}
Historically, rent controls in Canada were adopted at both provincial and federal levels. Canada introduced rent controls in 1944 during World War II as part of the National Housing Act. This "first-generation" of rent control froze nominal rents, and was inspired by the rent freeze that was adopted in European countries as a war-time measure during World War I, to protect veterans from evictions and rent increases. Whereas the rent freeze seemed justified during the war, maintaining it after the war started to negatively affect rental supply, labor mobility, and property maintenance \citep{arnott-1995}. Moreover, the "first-generation" of rent controls, largely studied in economics textbooks, was opposed by economists and landlords due to the widespread belief that it had resulted in significant deadweight losses. As a result, the Ontario provincial government repealed the "first-generation" of rent controls in 1949. The fact that the Ontario government decontrolled its rental markets is real-world evidence against the common belief that rent controls belong to a set of policies that are impossible to remove once adopted. 

In the wake of the oil crisis in the 1970's, the government of Ontario re-introduced rent controls in 1975 as part of the Residential Premises Rent Review Act. These new controls, classified as "second-generation" rent controls, differed significantly from "first-generation" rent controls because they had more flexible provisions allowing some mild nominal rent increases based on mandated guidelines. These controls also exempted luxury accommodation, government owned buildings, social housing, and private rental units built after January 1976. Under rent controls, only one rent increase was allowed each year in an amount that was not to exceed the rent increase set by government guidelines. Increases beyond the statutory percentage were granted in cases major increases in the cost of supplying rental housing, financial losses of landlords, or required capital expenditures. 

Since 1975, the Ontario rent control system has been evolving with two major changes. The first major change consisted of intermittent exemptions. The Residential Rent Regulation Act in 1985 removed exemptions for previously exempted units and newly built units. Under this new change, both new buildings and buildings constructed after 1976 were subject to rent controls. In 1992, rent control with exemptions for new construction was again restored by the Rent Control Act, which exempted all rental buildings built after November 1, 1991. However, in April 2017, the government of Ontario removed exemptions by extending rent controls to rental buildings built after November 1, 1991. At that time, all rental units in Ontario were under rent control. In 2018, the newly elected Ontario government enacted new legislation that only exempted buildings built, or first occupied, after November 15, 2018.

The second major change of the Ontario rent control system came in 1998, and was marked by tenancy deregulation as part of the Ontario Tenant Protection Act. According this Act, landlords, whose units were not exempted in Ontario, were allowed to increase rents by more than prescribed by the rent increase guidelines, but only when a signing a new contract with a new tenant. Otherwise, they were compelled to follow the rent increase guidelines. The 1998 Act was considered as a relaxation of the previous rent control because landlords were allowed to reset rents after a change of tenants. Unlike exemptions, no changes have been made to tenancy deregulation since that time. This type of rent control is sometimes referred to as tenancy rent control, partial rent control, or "third-generation" rent control \citep{arnott-2003}. 

This review highlights the evolution of the rent control system and its main characteristics that our model and experiments need to capture. It requires our model to capture exemptions, rent stabilization, and tenancy deregulation. It also requires to understand main changes from current rent control with exemptions to full rent control to full decontrol. These regime changes more or less have taken place. For example, in 2017 the government of Ontario removed exemptions, which brought all rental units under rent control; in 1949, the government of Ontario repealed rent control. Last it requires to understand both the aggregate and distributional outcomes. These elements are the focus of our modelling and experiments. 

\section{Model}\label{sec:model}

The model economy features overlapping generations of finitely-lived households, competitive rental companies, competitive housing construction companies, and a government. Households in each cohort are heterogeneous in labor market productivity, housing tenure, the type of rental market they are renting in (i.e., a rent-controlled market or a free market), years they have lived in the same rental unit, housing wealth, and financial wealth. Before retirement, households work in the labor market and after retirement they receive pensions from the government. Pensions depend on household average lifetime labor income. Homeowners pay proportional property taxes and maintenance costs each period, and mortgage debt is modelled with many of the features of long-term contracts. Competitive rental companies operate in a rent-controlled market and a free rental market. Housing units are supplied by competitive construction companies with convex costs. The government finances pensions to retired households, lump sum transfers, and public goods through taxes on income and properties. The economy is modelled as a small open economy and the interest rate, $r$, is exogenous \citep{sommeretal-2013,arslanetal-2015,sommersullivan-2017,kaplanetal-2019}. 

\subsection{Demographics}\label{subsection:demographics}

Households are born at age $j=1$, retire from the labor force at age $j=R$, and live to $J$ years. Households do not face mortality risk before age $J$, but after age $J$ they die with certainty. Households do not care about their children, and thus in this set up, optimally households die without assets. This means that newborn households begin their lives as assetless renters.

Similarly to \citet{sommersullivan-2017}, households have the following preferences:
\begin{equation*}
    u(c,h)=\frac{\left(c^{1-\gamma} h^{\gamma}\right)^{1-\sigma}}{1-\sigma},
\end{equation*}

\noindent where $c$ represents non-housing consumption, $h$ housing services consumption, $\gamma$ utility weight on the consumption of shelter services, and $\sigma$ household constant relative risk aversion. Households discount the future at a rate of $\beta$. %Non-housing consumption is taxed at a rate $\tau_{c}$. 

\subsection{Endowments}

We denote financial wealth by $a\in \mathcal{A}=\mathbb{R}_+$, tenure by $o \in \mathcal{O}=\{0,1\}$, the type of rental market by $f \in \mathcal{F}=\{0,1\}$, years lived in the same rental unit by $\iota \in \mathcal{I}=\mathbb{N}$, house size by $h\in \mathcal{H}=\mathbb{R}_{++}$, and mortgage debt by $m\in \mathcal{M}=\mathbb{R}_+$. Households are endowed with a stochastic labor market productivity, $e \in\mathcal{E} = \mathbb{R}_{++}$. The measure of age-$j$ households over endowments is $\Psi_{j}(a,o,f,\iota,h,m,e)$. We denote the state space of a household by  $\mathcal{S}=\mathcal{A}\times\mathcal{O}\times\mathcal{F}\times\mathcal{I}\times\mathcal{H}\times\mathcal{M}\times\mathcal{E}$. To simplify on notation, we mostly use $s$ to denote a household's state $(a,o,f,\iota,h,m,e)$.

Labor market productivity of a newborn is drawn from a conditional distribution \\ $\hat{F}(e^{\mbox{child}}|e^{\mbox{parents}})$, and it follows a Markov process $F(e'|e)$ until mandatory retirement. Labor market productivity also has a deterministic life-cycle component, $\zeta_j$, that increases with age; a household of age $j$ with labor market productivity $e$ supplies $e\zeta_j$ units of labor. Lastly, each household receives a lump sum transfer $T$ from the government each period.

\subsection{Labor and capital incomes}\label{subsection:income}

Working-age households with age below $R$ supply labor inelastically, and earn an exogenous wage, $W$, for each unit of labor they supply. Working-age households of age $j$ and of labor market productivity $e$ earn labor income of $y_{\ell,j}(e)=We\zeta_j(e)$. Retired households receive pensions, $y_{R}(e)$. Pensions depend on expected labor income conditional on labor productivity at retirement. Households that retire with a high level of productivity are more likely to have high productivity throughout their lives, and thus receive larger pensions. Households also earn capital income $y_{k}(a)=ar$ on their financial assets, $a$, where $r$ is the rate of return on financial assets. We define taxable income as

\begin{equation*}\label{eq:taxable_labor_income}
  y^{x}_{j}(e,a) =y_{\ell,j}(e)+\nu y_{k}(a),
\end{equation*}

\noindent where, $\nu$ is the fraction of capital income that is taxable. Then, following \citet{hsv-2020}, we define income taxes to be progressive as follows:

\begin{equation}\label{eq:income_tax}
  y^{\tau}_{j}(e,a) = y^{x}_{j}(e,a)-\tau_{\ell}y^{x}_{j}(e,a)^{\psi},
\end{equation}

\noindent where $\tau_{\ell}$ determines the average level of taxation in the economy and $\psi$ affects the degree of progressivity in the tax system.

\subsection{Housing}\label{subsection:housing_choice}

Households choose to rent or own a house, $o \in \mathcal{O}=\{0,1\}$, and house size, $h$. The smallest rental unit is $\underline{h}_{r}>0$, and the smallest owner-occupied house is $\underline{h}_{o}>0$ \citep{gervais,sommeretal-2013,sommersullivan-2017,kaplanetal-2019}. A household that chooses to own a house of size $h$ pays the purchase price $p h$. Buying and selling houses entail one-time proportional transaction costs, $\tau_{b} p h$ and $\tau_{s} p h$, respectively. As in \citet{fallis-1985}, renters can choose whether to rent in the rent-controlled market or the free rental market, where $f \in \{0,1\}$: $f=0$ stands for renting in the controlled rental market, while $f=1$ in the free rental market. 

While the uncontrolled rent is $p^{r}_{1}h$, the controlled rent depends on tenancy duration. In the controlled market, the longer the tenancy duration, the cheaper the real rent becomes. The controlled rent is a function of tenancy duration, $\frac{p^{r}_{0}}{(1+\kappa)^{\iota}} h$. Here, $p^{r}_{0}$ represents the rental price upon entering a controlled rental unit, $\kappa$ the rate at which rent control reduces real rents for each additional period a renter remains in the same unit, and $\iota$ the number of years a renter has lived in the same unit. 

In figure \ref{fig:diagram_rent_path} we show qualitatively illustrate the rents paid by renters as a function of the market in which they rent and their tenancy duration. The downward sloping red line shows rents paid by renters living in controlled units. When they enter a new rental unit, their rent is $p^{r}_{0}$, but as their tenure increase over time their rent declines. On the other hand, the horizontal red line shows rents paid by renters living in the free market. This rent remains constant regardless of their tenancy duration. In equilibrium, rents paid in the free market are lower than rents paid by renters in the controlled market early in their tenancy, but higher than the average rent in the controlled market. What explains these equilibrium outcomes? If initial rents paid by new tenants in the controlled market were equal or lower than the rents paid by renters in the free market, all renters would sort into the controlled market. To clear the markets, there has to be a mechanism in place to shift some renters to the free market. This mechanism is a premium paid initially in the controlled market. Renters who anticipate to remain only a few years in the same rental would sort into the free market to avoid paying the premium, whereas renters who expect to remain many years in the same rental would bite the bullet, pay the premium in the controlled market and enjoy lower average rents.

Each period, houses are taxed at a rate $\tau_{p}$ and depreciate in value at a rate $\delta_h$. Homeowners in the model receive an exogenous capital gain of $g_{h}ph$.\footnote{In Canada, capital gains from principal residences are fully non-taxable.} Rents in many highly-priced cities such as Toronto are lower than the ongoing costs of rental housing provision, ie property taxes, maintenance, opportunity cost of capital, and management. This implies that the relative price of homeownership is too high in the model, and so all households optimally choose to rent in equilibrium. To fix this, we introduce capital gains as in the standard user cost approach \citep{himmbelbergetal-2005}. We have also considered adding a utility advantage to homeownership to encourage households to own houses in equilibrium, but this approach does not fix the problem that with lower rent to price ratios, optimally no rental housing is supplied in equilibrium. Furthermore, we need an implausibly high utility advantage to match the homeownership rate. Therefore, we think that adding exogenous capital gains to housing is the best solution, considering the alternatives.\footnote{Alternatively, one could attempt to solve the model with aggregate uncertainty in house prices, but this approach is computationally infeasible considering all the other elements in our model.}

\subsection{Mortgages}\label{subsection:morgtages}

Mortgages are modelled with many features of long-term debt, including origination costs, prepayment penalties, and borrowing restrictions. We set mortgage rate to $r_{m}=r+\Delta r$, where $r$ is the return on financial wealth, and $\Delta r$ is the interest premium paid on mortgage debt \citep{sommeretal-2013,sommersullivan-2017,kaplanetal-2019}. Mortgages, $m$, are originated with a fixed cost, $\omega_{1}$. At the time of origination, the loan to value ratio (LTV) cannot exceed a certain fraction $\lambda_{m,1}(ph)$, defined as:

\begin{equation}\label{eq:LTV}
m \leq \lambda_{m,1}(ph) p h.
\end{equation}

\noindent $\lambda_{m,1}(ph)<1$ and depends on the value of the house at purchase. Mortgage debt at origination is also limited by a gross debt servicing (GDS) ratio: 

\begin{equation}\label{eq:GDS}
\frac{\tau_{p}ph+(r_{m}+\upsilon_{3})m}{y_{\ell,j}(e)+y_{k}(a)} \leq \lambda_{m,2}\left(\frac{m}{ph}\right).
\end{equation}

\noindent $\upsilon_{3}$ is the average fraction of debt a household pays over the life of a mortgage, and $\lambda_{m,2}\left(\frac{m}{ph}\right)$ is the maximum GDS ratio allowed, which depends on the LTV ratio. The GDS constraint (\ref{eq:GDS}) implies that GDS--including property taxes, interest payments, and the average fraction of the principal paid--cannot exceed a fraction, $\lambda_{m,2}\left(\frac{m}{ph}\right)$, of their pre-tax income from all sources. Each period, owners with a mortgage, $m>0$, must pay a fraction $\upsilon_{1}$ of the total mortgage debt, and are allowed to pay up to a fraction $\upsilon_{2}>\upsilon_{1}$ of the debt without a penalty. Any amount paid above $\upsilon_{2}m$ is subject to a prepayment penalty as follows: 

\begin{equation}\label{eq:prepayment_penalty}
\tau_{pp}=\omega_{2}\mbox{max}\left\{(1-\upsilon_{2})m-m',0\right\}
\end{equation}

\noindent If an owner does not pay at least $\upsilon_{1}m$, so that $m'>(1-\upsilon_{1})m$, the homeowner is considered to have originated a new mortgage, and thus must pay an origination cost, $\omega_{1}$, as well as satisfy borrowing constraints (\ref{eq:LTV})-(\ref{eq:GDS}). Once mortgage debt is below a certain threshold relative to the value of the house, $m \leq \upsilon_{4}ph$, the homeowner must completely repay the remainder of its debt by the end of the period, and is not subject to prepayment penalties. Mortgage debt is non-portable and retired households cannot originate new debt, but can carry on existing debt as detailed above. Because households in Canada are not allowed to deduct their mortgage interest payments or property taxes from their taxable income, they cannot do so in the model either.

\subsection{Dynamic program}\label{subsection:household_optimization}
Given its state variables, a household chooses how much to save ($a'$), whether to own or rent ($o' \in \left\{0,1\right\}$), whether to rent in the rent-controlled market or the free rental market ($f' \in \left\{0,1\right\}$), and whether to remain in their rental unit one period longer ($\iota'$). The household also chooses how much housing to consume ($h$), and how much mortgage debt to take on ($m'$). The dynamic program that represents this problem in a stationary equilibrium is

\begin{align}
  \label{eq:vf}
  V_{j}(s)=\max_{c,a',o',f',\iota',h',m'}\ \biggl\{u\left(c,h\right)+
  \beta \mathbb{E}\biggl[V_{j+1}(s'|e)\biggl]\biggr\}
\end{align}
subject to
\begin{flalign}
  \begin{split} \label{eq:bc}
  & c+a'+(1-o')\left[(1-f')\frac{p^{r}_{0}}{(1+\kappa)^{\iota'-1}}h'+f'p^{r}_{1}h'\right]+o'(ph'-m')+o\left(ph\left(\tau_{p}+\delta_{h}\right)+r_{m}m\right) + \\ & \quad \indicator{(o'=1\land o=0)\vee\left(o'=1 \land o=1 \land h' \neq h\right)}\tau_{b}ph' +\indicator{(o'=0\land o=1)\vee\left(o'=1 \land o=1 \land h'\neq h\right)}\tau_{s}ph+\\ & \quad \omega_{1}\indicator{\left(m'>(1-\upsilon_{1})m\right)\vee \left((m'>0) \land (m\leq\upsilon_{4}ph)\right)}+\tau_{pp}\indicator{(m'<(1-\upsilon_{2})m)\land \left(m>\upsilon_{4}ph\right)} = \\ & \quad  y_{\ell,j}(e)+y_{k}(a)-y^{\tau}_{j}(e,a)+a+o(ph-m)+og_{h}ph+ T   
  \end{split} \\
  & c \geq 0 \label{eq:consumption}\\
  & a' \geq 0 \label{eq:financial_wealth}\\
  & o' \in \{0,1\} \label{eq:tenure}\\
  & f' \in \{0,1\} \hspace{0.2cm} \mbox{if} \hspace{0.2cm} o'=0 \label{eq:rental_market_type} \\
  & h' \geq \underline{h}_{r} \hspace{0.2cm} \mbox{if} \hspace{0.2cm} o'=0 \label{eq:min_housing_renter}\\
  & h' \geq \underline{h}_{o} \hspace{0.2cm} \mbox{if} \hspace{0.2cm} o'=1  \label{eq:min_housing_owner}\\
  & \lambda_{m,1}(ph') p h'\geq m'\geq 0 \hspace{0.2cm} \mbox{if} \hspace{0.2cm} \indicator{\left(m'>(1-\upsilon_{1})m\right)\vee \left((m'>0) \land (m>\upsilon_{4}ph)\right)}=1 \hspace{0.2cm} \mbox{and} \hspace{0.2cm} j\leq R\label{eq:LTV_1}\\
  & \frac{\tau_{p}ph'+(r_{m}+\upsilon_{3})m'}{y_{\ell,j}(e)+y_{k}(a)} \leq \lambda_{m,2}\left(\frac{m}{ph}\right) \hspace{0.2cm} \mbox{if} \hspace{0.2cm} \indicator{\left(m'>(1-\upsilon_{1})m\right)\vee \left((m'>0) \land (m\leq\upsilon_{4}ph)\right)}=1 \hspace{0.2cm} \mbox{and} \hspace{0.2cm} j\leq R\label{eq:GDS_1} \\
  &\tau_{pp}=\omega_{2}\mbox{max}\left\{(1-\upsilon_{2})m-m',0\right\} \hspace{0.2cm} \mbox{if}  \hspace{0.2cm}\indicator{(m'<(1-\upsilon_{2})m)\land \left(m>\upsilon_{4}ph\right)}=1\label{eq:prepayment_penalty_1}\\
  & m'=0 \hspace{0.2cm} \mbox{if} \hspace{0.2cm} o'=0. \label{eq:mortgage_debt_renter}
\end{flalign}

\noindent Household utility includes two components: utility from housing and non-housing consumption; expected value based on its choices discounted by discount factor, $\beta$.\footnote{We include labor market productivity shocks for retired households as well as non-retired households in our formulation of the dynamic program for notational brevity. Writing the dynamic program this way allows us to use a single equation to represent both retired and non-retired households' problems.} 

Equation \ref{eq:bc} represents household budget constraint. The left-hand side contains standard expenditures: consumption, $c$; saving, $a'$ ; the cost of renting a unit in the rent controlled market as a function of the years lived in the same rental unit, $\frac{p^{r}_{0}}{(1+\kappa)^{\iota'-1}}h'$ ;the cost of renting in the free rental market, $p^{r}_{1}h'$; and the purchase price of a house using a mortgage for homeowners, $(ph'-m')$. The left-hand side also includes costs to homeownership: property taxes $ph \tau_{p}$, maintenance costs $ph \delta_{h}$, and mortgage interest costs for a household starting the period as a homeowner $r_{m}m$. The other terms include: buying costs for a household starting the period as a renter and that chooses to become a homeowner or for a household starting the period as a homeowner that decides to purchase a different property, $\tau_{b}ph'$; selling costs for a household starting as an owner and that decides to move to a different unit, $\tau_{s}ph$; origination costs for owners who take on new mortgage, or do not pay at the minimum payment required on their existing debt, $\omega_{1}$; and prepayment penalty for owners with a mortgage who pay more than $\upsilon_{2}$ of their debt and whose debt is above a fraction $\upsilon_{4}$ of their house value, $\tau_{pp}$. 

The right-hand side of the budget constraint contains resources available to spend: after-tax labor and capital incomes, $y_{\ell,j}(e)+y_{k}(a-y^{\tau}_{j}(e,a)$; non-housing financial wealth, $a$; housing equity for a household that begins the period as an owner, $(ph-m)$; capital gains from housing for owners, $g_{h}ph$; and a lump sum transfer, $T$. 

Constraints \ref{eq:consumption} and \ref{eq:financial_wealth} restrict households to have positive consumption and savings. Constraint \ref{eq:tenure} specifies the tenure choice of the value $1$ for owners, or $0$ for renters. Constraint \ref{eq:rental_market_type} indicates the choice of renters between renting in the free rental market, $f=1$, and renting in the rent-controlled market, $f'=0$. Constraints \ref{eq:min_housing_renter} and \ref{eq:min_housing_owner} require renters and owners to consume at least $\underline{h}_{r}$ and $\underline{h}_{o}$ units of housing, respectively; constraints \ref{eq:LTV_1} and \ref{eq:GDS_1} require the origination of mortgage debt to satisfy the LTV and GDS ratio conditions. Constraint \ref{eq:prepayment_penalty_1} details the prepayment penalty for owners that pay more of their debt than allowed. Constraint \ref{eq:mortgage_debt_renter} says that renters must have no mortgage debt. 

When solving for the problem of households, we denote the policy functions for consumption, savings, tenure, rental market choice, years living in the same unit, housing, and mortgage debt by $c_{j}(s)$, $a'_{j}(s)$, $o'_{j}(s)$, $f'_{j}(s)$, $\iota'_{j}(s)$, $h'_{j}(s)$, and $m'_{j}(s)$ respectively.

\subsection{Rental company}

We follow the rental housing supply theory in \citet{fallis-1985} and embed it \citet{garrigaeetal-2009}'s rental housing supply structure. For each unit, $x_{i}=1$, of rental housing, the rental company obtains average rent, $\overline{p}^{r}_{i}$, pays property taxes, $\tau_{p}h$, and incurs depreciation costs, $\delta_{h}h$, where $i=0$ represents the controlled rental market, and $i=1$ denotes the free rental market. Like homeowners, the rental company also receives a non-taxable capital gain, $g_{h}p$, each period on each unit. Because rents do not change with tenancy duration in the free rental market, $\overline{p}^{r}_{1}=p^{r}_{1}$. 

To supply a total of $X'=\sum^{1}_{i=0} X'_{i}$ units of rental housing, the rental company also incurs management costs equal to $c_{m}X'^{\theta}$, where $c_{m}$ determines the average cost to manage the stock of rental housing $X'$ and $\theta$ determines how quickly the marginal cost of supplying rental housing rises with the amount of rental housing supplied. 

The rental company must supply $\mbox{min}\{\underline{X}_{0},X'\}$ in the controlled market, with $\underline{X}_{0} \geq 0$. This means that if the total supply of rental housing is $X'\leq \underline{X}_{0}$, the market is completely under control, and the supply of rental housing is determined by the average rent in the controlled market, $\overline{p}^{r}_{0}$. However, if $X'>\underline{X}_{0}$, the total supply of rental housing, $X'$, is a function of the free market rent $\overline{p}^{r}_{1}$, where $\underline{X}_{0}$ is the supply of rent-controlled units and $X'-\underline{X}_{0}$ is the supply in the free market.

The rental company, without friction, buys houses at the house price, $p$, and each period decides how much housing to supply in the rent-controlled market, $X'_{0}$, and how much housing to supply in the free rental market $X'_{1}$. Thus, the rental company has the following problem:

%Because the rent-to-price ratio in Toronto is quite low, it does not cover the flow of operating costs, which would imply a net loss to the rental company. To fix that, we add another parameter, $g_{p}$, which exogenously accounts for the real growth in house prices in Toronto, and thus guarantees that the rental company has an incentive to supply rental housing.

\begin{align}\label{eq:rental_company_problem}
V(X)=\underset{X'_{0},X'_{1}} {\text{$\mbox{max} $}}\left\{\sum_{i=0}^{i} \overline{p}^{r}_{i}X'_{i}-c_{m}X'^{\theta}-p X'+p(1-\delta_{h}-\tau_{p}+g_{p})X+\frac{1}{1+r}V(X')\right\}.
\end{align}
subject to
\begin{align}
& X'_{0}=\mbox{min}\{\underline{X}_{0},X'\} \label{eq:supply_restriction} \\
& X'_{1} \geq 0 \label{eq:supply_positive} \\
& X'=\sum^{1}_{i=0} X'_{i} \label{eq:total_supply} 
\end{align}

\noindent Equation \ref{eq:supply_restriction} says that the rental supply in the controlled market is $\mbox{min}\{\underline{X}_{0},X'\}$. Equation \ref{eq:supply_positive} says that rental housing supply in the free market cannot be negative. Equation \ref{eq:total_supply} says that total rental housing supply is the sum of the rental supply in the rent-controlled and free markets. Since newly built rental units are exempted from rent control\footnote{This is also how rent control is done in Toronto.}, the marginal revenue of the rental company is determined by the free market rent whenever $X'>\underline{X}_{0}$. Thus, the total rental supply is given by the following inverse supply equation:

\begin{align}
\overline{p}^{r}_{1}(X')=c_{m}\theta X'^{\theta-1}+p-\left\{\frac{1-\delta_{h}-\tau_p+g_{p}}{1+r}\right\}p.\label{eq:rental_price}
\end{align}

\noindent In equilibrium, when $X'\geq \underline_{X}_{0}$, then $X'_{0}=\underline_{X}_{0}$ and $X'_{1}=X'-\underline{X}_{0}$. If $X'<\underline_{X}_{0}$ or if regulation requires all rental units to be under rent control, such that $\underline{X}_{0}=X'$ all the time, then $\overline{p}^{r}_{0}$ replaces $\overline{p}^{r}_{1}$ in equation \ref{eq:rental_price}, $X_{0}=X'$, and $X_{1}=0$.

\subsection{Housing construction company}

Following \citet{rotberg-2022}, houses are supplied endogenously by a construction sector. The law of motion for housing is $H'=(1-\delta_h)H + I$, where $I$ denotes the construction sector's investment in new housing stock. For its production of $I$ new units, the construction company receives payments equal to $pI$, but incurs an increasing marginal cost, as captured by the following production technology:

\begin{equation*}
TC(I)=c_{1}\left(\frac{1}{\frac{1}{\epsilon}+1}\right)I^{\frac{1}{\epsilon}+1}\label{eq:total_cost_structure},
\end{equation*}

\noindent where $c_1$ scales up or down the cost of producing new housing and $\epsilon$ represents the price elasticity of housing investment. The construction company chooses investment to maximize profits:

\begin{equation*}
\label{eq:construction_company_profits}
\underset{I} {\text{$\mbox{max} $}}{\left\{p I-TC(I)\right\}},
\end{equation*}

\noindent The solution to this problem gives the following relationship between housing investment and the price of housing:

\begin{equation}
\label{eq:construction_company_first_order_condition}
I=\left(\frac{p}{c_{1}}\right)^{\epsilon}.
\end{equation}

\subsection{Aggregation}\label{subsection:aggregation}

The government's budget is balanced in all our experiments and in the benchmark. The government's budget constraint is as follows.
\begin{align}\label{eq:gbc}
  G + \sum_{j=R}^J \int_\mathcal{S} y_{R}(s) \mathrm{d}\Psi(s) + \sum_{j=1}^J \int_\mathcal{S} T \mathrm{d}\Psi(s) = \sum_{j=1}^J \int_\mathcal{S}y^{\tau}_{j}(s) \mathrm{d}\Psi(s)+\tau_p p H,
\end{align}

\noindent The budget constraint states that the government's expenditures on public goods, pension payments, and lump sum transfers must equal to tax revenues from income and property. When calibrating the model to the current tax code in Toronto, we set the lump sum transfer $T$ to 0 and choose the consumption level of public goods, $G$, to ensure that this constraint holds. In our experiments, we keep $G$ constant, and adjust $T$ to clear the budget constraint. Note that $S=A\times O \times F \times I \times H \times M \times E$ denotes a typical subset of the state space $\mathcal{S}$, $\mathcal{A}$, $\mathcal{O}$, $\mathcal{F}$ $\mathcal{I}$, $\mathcal{H}$, $\mathcal{M}$, and $\mathcal{E}$, respectively.

The overall demand for housing equals the overall supply of housing:

\begin{equation}
\label{eq:housing_clearing_condition}
  \sum_{j=1}^J \int_\mathcal{S} h_{j}(s) \mathrm{d}\Psi(s) = H',
\end{equation}

\noindent and the overall demand for rental housing in each rental market must equal the overall supply of rental housing in each market:

\begin{equation}
\label{eq:rental_housing_clearing_condition}
  \sum_{j=1}^J \int_\mathcal{S} (1-o'_{j}(s))\indicator{f'_{j}(s)=i}h_{j}(s) \mathrm{d}\Psi(s) = X'_{i} \quad i \in \{0,1\}.
\end{equation}

The distributions of households evolve according to the laws of motion,
\begin{align}
\Psi_{j+1}(S) &=\int_{\mathcal{S}} F(E,e)M_{j}(s)\mathrm{d}\Psi_{j}(s),\ j=1,\ldots,J-1\label{eq:lom1}
\end{align}
\begin{align}
\Psi_{1}(a=0,o=0,f=0,\iota=0,h=\underline{h}_{r},m=0,e) &=\int_{\mathcal{S}}\tilde{F}(E,e) \mathrm{d}\Psi_{J}(s),\label{eq:lom2}
\end{align}

\noindent where,

\begin{align}
    M_{j}(s) = \indicator{[a'_{j}(s)\in A] \land [o'_{j}(s)\in O] \land [f'_{j}(s)\in F] \land [\iota'_{j}(s)\in I] \land [h'_{j}(s) \in H] \land [m'_{j}(s) \in M]}.\notag
\end{align}

\subsection{Equilibrium}
A stationary equilibrium in the model is a sequence of aggregate prices and quantities, $\left\{H',I,\{X'_{i},\overline{p}^{r}_{i}\}^{1}_{i=0},p\right\}$, a sequence of value and policy functions, $\biggl\{V_{j}(\cdot),c_{j}(\cdot),a'_{j}(\cdot),o'_{j}(\cdot),f'_{j}(\cdot)\iota'_{j}(\cdot),h'_{j}(\cdot),m'_{j}(\cdot)\biggl\}_{j=1}^J$, and a sequence of measures, $\left\{\Psi_{j}(\cdot)\right\}_{j=1}^J$, that:
\begin{enumerate}
%\item all nominal prices rise at the rate of $i$, but real prices remain fixed;
\item Solve the household's dynamic problem (\ref{eq:vf}) subject to constraints (\ref{eq:bc})-(\ref{eq:mortgage_debt_renter});
\item Solve the rental company's problem (\ref{eq:rental_company_problem}), subject to (\ref{eq:supply_restriction})-(\ref{eq:total_supply});
\item Satisfy the construction company's first-order condition, (\ref{eq:construction_company_first_order_condition});
\item Satisfy the government's budget constraint, (\ref{eq:gbc});
\item Clear the housing market,
(\ref{eq:housing_clearing_condition});
\item Clear the rental housing markets,
(\ref{eq:rental_housing_clearing_condition});
%\item satisfy the average declared house price as described in
%(\ref{eq:average_declared_price});
\item And satisfy the laws of motion for the distribution of households, (\ref{eq:lom1})--(\ref{eq:lom2}).
\end{enumerate}

\section{Calibration and model performance}\label{sec:calibration}

Calibration follows two steps. First, we calibrate certain parameters either with estimates from the literature or from data in Toronto. Second, we internally calibrate the values of the remaining parameters so that the model generates certain statistics as observed in the data. Table \ref{tab:calibration} lists the parameter values. 

\subsection{Assigned parameter values}

\subsubsection{Demographics, preferences, and labor market productivity}
A model period corresponds to 1 year in the data. Households start their lives at age 26, retire at age 66, and live to age 85, implying that lifespan $J=60$ and retirement age $R=41$. We set households' relative risk aversion $\sigma=2$. The discount factor, $\beta$, and the share of housing services in utility, $\gamma$, are determined in the second stage of our calibration procedure.

Labor market productivity of households and inter-generational labor market productivity follow AR(1) processes,
\begin{equation*}
  \log e' = \rho_{e} \log e + \epsilon_{e},\quad \epsilon_{e}\sim N(0,\sigma_{e}),
\end{equation*}
\begin{equation*}
\log e^{\text{child}} = \hat{\rho}_{e} \log e^{\text{parent}} + \hat{\epsilon}_{e},\quad \hat{\epsilon}_{e}\sim N(0,\hat{\sigma}_{e}).
\end{equation*}
We set the intra-generational persistence, $\rho_{e}=0.852$, and the variance, $\sigma_{e}=0.376$, based on our own estimates using linked household-level micro-data of the 2011 and 2016 waves of the Census.\footnote{Because we use Census data from 2011 and 2016, our estimates are for a period of 5 years. Therefore, we estimate $\rho_{e,5}$ and $\sigma_{e,5}$. Using backward substitution one can show that $\rho_{e}$ is $\rho_{e,5}$ raised to the power of 1/5 and that $\sigma_{e}$ is $\sigma_{e,5}$ divided by square root of 5.} The inter-generational persistence, $\hat{\rho}_{e}$ is set to $0.199$ and the variance, $\hat{\sigma_{e}}=1.09$, based on estimates for Canada from \citet{corak-heisz-1999}. Using 2016 Census data, we set the deterministic life-cycle component of labor market productivity to 
\begin{equation*}
\zeta_j = \begin{cases}
1+1.48 \frac{j}{R-10} &\text{if $j\leq R-10$}\\
\zeta_{R-10}-\left(j-\left[R-10\right]\right) (1.48-1.04) \frac{j}{10} &\text{if $R \geq j>R-10$}.
\end{cases}
\end{equation*}

\noindent This formulation reflects the data: real income is 148\% higher for households aged 55 than those aged 26 and real income of households aged 65 is 104\% higher than real income of households aged 26. We approximate pension benefits, $\{y_{R}(e)\}_{e\in \mathcal{E}}$, using the estimates of the replacement rate of gross individual earnings as a function of average lifetime income relative to the average income in the economy \citep{OECD-pension}.

\subsubsection{Taxes, interest rate, and housing}

The portion of capital income that is taxable, $\nu$, is set to 50\% as in Canada's tax code. Using micro-data for Toronto from the 2016 Census, we estimate $\psi$ in equation (\ref{eq:income_tax}), and set $\psi=0.928$ to our point estimate of $\psi$. The interest rate, $r$, is set to 2.77\% as estimated for Canada by \citet{bruneau-etal-2018}. 

%We set the consumption tax rate to 13\%, which is the Harmonized Sales Tax (HST) in Ontario.

We set housing costs and transaction costs using a variety of data sources. We set the rate at which real rents decrease for renters remaining in the same unit in the rent-controlled market, $\kappa$, to equal to 2.70\%. We estimate this using Census data from 2001-2016 on the share of the rent-controlled market, and Rental Market Survey (RMS) data for 1992-2021 on the turnover rate of renters and average real growth of average rents. We set capital gains from housing, $g_{p}$, to 4.46\% based on estimates from Multiple Listing Service (MLS) data from 2006-2016, which show that real house prices in Toronto rose by about 4.46\% annually. We set the property tax rate, $\tau_p$, to 0.67\% based on 2016 data from over 20 municipalities in Toronto, and the housing depreciation rate to 2\% as estimated by \citet{hardingetal-2007}. We set buying costs, $\tau_{b}$, to 2.5\% of the purchase price, and selling costs, $\tau_{s}$, to 7\% based on estimates by \citet{gm-2003}. We choose $\epsilon$ to equal 0.35 to align with the \citet{cmhc-report-2018} estimate of the elasticity of the housing supply in Toronto. Lastly, we set the free market rent, $p^{r}_{1}$, to 5.85\% of the equilibrium house price to be consistent with our estimate of the rent-to-price ratio in both the rent-controlled and free markets, which we obtain using data for 2016 from the Canadian Real Estate Association (CREA).

\subsubsection{Mortgages}

Based on banking regulations in Canada, we set the maximum ratio of mortgage debt to house value, $\lambda_{m,1}(ph)$, as follows:

\begin{equation*}
\lambda_{m,1}(ph) = \begin{cases}
0.8 &\text{if $ph\geq \$1M$}\\
0.9 &\text{if $\$500K\leq ph< \$1M$}\\
0.95 &\text{if $ph < \$500K$}.
\end{cases}
\end{equation*}

\noindent The above conditions specify that: homeowners with a house that has a model-equivalent value above \$1M can borrow only up to 80\% of its value; owners purchasing a home with a model-equivalent value between \$500,000-\$1M can borrow up to 90\% of the house value; and owners with a house value below \$500,000 can get a loan of up to 95\% of the house value. 

The maximum GDS ratio, $\lambda_{m,2}\left(\frac{m}{ph}\right)$ takes the following form:

\begin{equation*}
\lambda_{m,2}\left(\frac{m}{ph}\right) = \begin{cases}
0.39 &\text{if $\frac{m}{ph}>0.8$}\\
0.47 &\text{if $\frac{m}{ph}\leq 0.8$}.
\end{cases}
\end{equation*}

\noindent This equation says that for mortgages with a LTV ratio above 80\% the maximum GDS ratio is 39\%, and that the maximum GDS ratio for mortgages with a LTV below 80\% is 47\%, which we have learned from big lenders in Canada that this is oftentimes the limit in cases of lower-risk mortgage loans. 

We calibrate mortgage contracts with great details. We set the fixed mortgage origination cost, $\omega_{1}$, to 2.28\% of the average household income in Toronto. The mortgage origination cost includes the cost of appraisals ($\approx \$500$) and lawyer fees ($\approx \$1,200$). We set the prepayment penalty, $\omega_{2}$, to 3 months of interests, a common practice in variable-rate mortgage contracts in Canada. We set the fraction a homeowner must pay each period, $\upsilon_{1}$, to avoid the cost of originating a new loan to 2.9\%, a choice which we explain shortly below. We set the threshold for prepayment penalty, $\upsilon_{2}$, to $\upsilon_{1}$ plus 15\% of the total debt, 15\% of the required payment $\upsilon_{1}$, and 15\% of interest payments, which adds up to 18.8\%.\footnote{Standard mortgage contracts in Canada normally allow prepayments of 15\% of the principal and regular payments without penalty.} We set the average fraction of a mortgage paid each year, $\upsilon_{3}$, which is used in inequality (\ref{eq:GDS}), to 3.33\% to reflect a 30-year amortization period, a common amortization period in the mortgage industry in Canada. We also set the debt to house value ratio below which a homeowner must fully settle their debt, $\upsilon_{4}$, to 40.3\% assuming that the debt at origination was 95\% of the house value.\footnote{About 40\% of a mortgage loan amortized over 30 years is paid in the last 5 years.} The combination of $\upsilon_{1}$ and $\upsilon_{4}$ minimizes the sum of squared distances between the outstanding debt of a household making regular payments in the model each period, and the outstanding debt of a 30-year mortgage of a household making regular payments in reality. Finally, we set the interest rate premium on mortgage debt, $\Delta r_{m}$, to 0.19\%, based on the average spread in 2002-2022 between the 5-year variable rate on mortgages and the return on a 10-year government bond.

\subsection{Internally calibrated parameters}

We now turn to internally calibrating our model. After assigning the parameter values listed above, we still need to calibrate seven parameters: the rate of time preference, $\beta$; the utility weight on the consumption of housing services, $\gamma$; the smallest dwelling for purchase, $\underline{h}_{o}$; the average level of income taxation, $\tau_{\ell}$; the convexity of managing rental housing, $\theta$; the rent a tenant pays in the first year upon entering a rental unit under rent control, $p^{r}_{0}$; and the smallest housing unit, $\underline{h}_{r}$.\footnote{The scalar for construction costs, $c_{1}$, is chosen such that housing supply is exactly equal to housing demand in the benchmark equilibrium. In counterfactual simulations $c_{1}$ is held fixed at its benchmark level. The rental management cost, $c_{m}$, is chosen to equate the supply of rental housing with the demand for rental housing in the benchmark. In counterfactual simulations it is held fixed at its benchmark level.}

We jointly calibrate all these parameters to match several moments in Toronto data. The percentage of aggregate mortgage debt in aggregate gross housing wealth, which is 24.3\% according to the 2016 wave of the Survey of Financial Security (SFS). The percentage of aggregate renters' income that renters spend on rent, which is 24.3\% and comes from 2016 Census. The homeownership rate, which is equal to 66.5\% according to the 2016 Census. The average labor income tax rate, which is 19.9\% and comes the 2016 Census. The elasticity of rental housing supply, which we estimate to be 0.2 using data from the 2006 and 2011 Census and MLS. The free rental market share in the total rental market, which is 12.4\% and comes from the 1991-2016 waves of the Census. The percentage of renters spending over 30\% of their gross income on rent, which comes from the 2016 Census and is equal to 37.5\%.

In general, the rate of time preference, $\beta$, controls the aggregate mortgage debt to aggregate gross housing value ratio; the utility weight on housing, $\gamma$, determines mostly the fraction of aggregate renters' income renters spend on rent; the smallest dwelling for purchase, $\underline{h}_{o}$, mainly affects the homeownership rate; the average income taxation, $\tau_{\ell}$, determines the average income tax rate; the convexity of managing rental housing, $\theta$, determines the elasticity of rental housing supply; the initial rent to pay for entering a controlled unit, $p^{r}_{0}$, determines the share of the free market; and the smallest housing unit, $\underline{h}_{r}$, controls the percentage of renters spending over 30\% of their gross income on rent.

\subsection{Model performance}\label{subsection:non_targeted_statistics}

We evaluate our model performance with aggregate statistics, life-cycle statistics, and tenure distribution. We start with aggregate statistics. In table \ref{tab:aggregate_non_targreted_statistics} we show the performance of our model on aggregate dimensions not targeted in our calibration. The model produces a reasonable although somewhat higher average owned-house value to average household income ratio relative to the data, and predicts a slightly lower percentage of homeowners with a mortgage compared to the data. We do not have an estimate for the ratio of the average value of rental units to the average value of owner-occupied units in Toronto. However, we can compare our model's estimate to \citet{chatterjee-eyigungor-2015}'s estimate for the U.S., and we find that our model is somewhat lower than their estimate. Using RMS data from 2016 through 2021, we estimate that the average tenure duration in the same rental unit ranges between 6.3 and 11.5 years, whereas our model produces an estimate that is equal to the lower bound of this estimate. Our model produces similar rent gap between controlled and uncontrolled rental units of 13.7\% as in \citet{zhang-gou} using RMS data in Toronto.

Moreover, similar to data, in our model the average income and house size of renters in the free rental market are 36\% and 30\% higher than in the rent-controlled market, respectively. The percentage of renters who spend over 50\% of their income on rent in our model is 9.7\%, which is quite similar to the estimate from the 2016 Census. The wealth distribution in our model is less concentrated than in data because our model cannot generate the high concentration of net wealth at the very top of the wealth distribution. This is because our model does not feature much rate-of-return heterogeneity\footnote{In the model homeowners receive an exogenous 4.46\% return on their house, and households with financial assets receive a 2.27\% return on those assets.}, which \citet{fagereng} show is important for generating a high concentration of wealth at the top of the wealth distribution. Alternatively, our model also does not feature heterogeneous discount factors \citep{ks-1998} or Michael Jordan income shocks \citep{mj-shocks}, which could also explain its inability to generate the high concentration of wealth at the top of the wealth distribution.

We also show the performance of our model on life-cycle dimensions As for life-cycle in figure \ref{fig:life_cycle}. As panel (a) shows, the homeownership rate in our model over households' life-cycle traces the data quite closely. However, as shown in panel (b), households in our model have somewhat higher LTV ratios compared to the data for young households, but somewhat lower for old households. 

Lastly, we show tenure by labor productivity in table \ref{tab:tenure_by_labor_productivity}, where $e_{1}$ represents the lowest labor productivity and $e_{7}$ the highest. We break down tenure into four categories: renters in the rent-controlled market, renters in the free market, homeowners with positive mortgage debt, and homeowners without any mortgage debt. The calibrated model performs exactly as one would expect. Households with lower $e$'s tend to be renters, while households with higher $e$'s tend to be homeowners without a mortgage. As a result of this, the percentage of households who own a home with a mortgage initially rises with $e$ and then declines. As well, households with lower $e$'s tend to rent in the rent-controlled market since they expect to remain renters for longer, whereas households with higher $e$'s are more likely to rent in the free market since they expect to become owners more quickly. 

\section{Experiments}\label{section:experiments}

\paragraph{Experiments' overview} Using our calibrated model, we conduct four experiments. In the first experiment, we compare the benchmark with exemptions and tenancy deregulation with the long-run equilibrium in which there is no rent control at all, which we also refer to as the full decontrol equilibrium. In the second experiment, we compare the benchmark to the long-run equilibrium in which all rental housing is controlled, which we also refer to as the full control equilibrium. In the third experiment, we search for the degree of rent control which produces the highest welfare in the long-run equilibrium. In our fourth experiment, we compute the transitional dynamics from the benchmark to the long-run equilibrium with full decontrol. In this experiment, households are hit by an unexpected shock, where the government fully decontrols rental markets permanently. In all our experiments we keep the tax code fixed and balance the government's budget with lump sum transfers. 

\paragraph{Welfare measure} Our main measure of welfare is the same as that used by \citet{conesa-et-al} and \citet{gkkoc}, which measures the expected ex-ante utility of a newborn household. It is equivalent to giving all newborn households a fixed proportional non-housing consumption transfer. Formally, the non-housing consumption-equivalent welfare effect due to a change in policy after $t$ periods is given by

\begin{equation*}
  CE_{t} = \left[\frac{\sum_{j=1}^{k}\int_{\mathcal{S}}V^{\dagger}_{j,t}(s)\ \mathrm{d}\Psi^{\dagger}_{j,t}(s)}{\sum_{j=1}^{k}\int_{\mathcal{S}}V^*_{j,t}(s)\ \mathrm{d}\Psi^*_{j,t}(s)}\right]^{\frac{1}{(1-\gamma)(1-\sigma)}}-1,
\end{equation*}

\noindent where daggers denote equilibrium objects in the counterfactual equilibrium and stars denote equilibrium objects in the benchmark. For our main welfare analysis, which considers newborn households only, we set $k=1$. We also present average welfare gains or losses across all households by setting $k=J$.

\paragraph{Measures of rent} Our model features several different measures of rent. Here we a provide a list to clarify each of these measures. The list is as follows: 

\begin{enumerate}
    \item \underline{Free market rent}: it is the equilibrium rent in the free market.
    \item \underline{Average controlled rent}: it is the equilibrium average rent in the controlled market.
    \item \underline{Rent of new tenants}: it is the equilibrium rent new tenants pay when moving into a new rental unit. In the controlled market, new tenants pay more than tenants who have been under rent control for a while. In the free market, new tenants pay the same rent as tenants that have been living in the same rental unit in the free market for a while.
    \item \underline{Average rent}: it is the weighted average of equilibrium rents in both the rent-controlled and free rental markets.
    \item \underline{Marginal rent}: it is the rent the rental company would receive if it supplied an additional rental unit. When free markets operate, the marginal rent is the equilibrium free market rent, which is higher than the equilibrium average rent in the controlled market. When all rental units are under control, the marginal rent is the equilibrium average rent in the controlled market.
\end{enumerate}

\subsection{Benchmark vs. full rent decontrol}
Our first experiment compares the benchmark equilibrium with the long-run equilibrium in which rental markets are completely decontrolled. We first present aggregate results, then distributional ones, and lastly, break down welfare changes into three effects: rent effect, house prices effect, and lump sum transfer effect.

\subsubsection{Aggregate results} The second column of table \ref{tab:aggregate_effects} presents changes in aggregate variables in the decontrolled equilibrium relative to the benchmark. Fully decontrolling the rental market would increase house prices, homeownership, and average wealth of homeowners. It would reduce the value of staying longer in the same rental unit, and thus encourage renters to save more to buy a house and increase the average wealth of renters by 1.09\%. As a result homeownership would rise by 0.65p.p.. Moreover, the decline in the value of remaining in the same rental unit would encourage renters to move more frequently and reduce their stay in the same rental unit by 54.6\% and demand 1.3\% larger units over time as their incomes grow. The two effects combined would raise the demand for housing, and thus increase house prices by 0.8\%. The increase in house prices would raise average wealth of homeowners by 0.56\%. as well as property tax revenues, and thus enable the government to provide lump sum transfers equal to 0.03\% of the average labor income. 

Decontrolling would also affect rental markets in several ways. Decontrolling would increase the average rent of previously-controlled units in the benchmark by 5.2\%. However, decontrolling would reduce the average rent of previously-uncontrolled units in the benchmark by 3.3\% because it would raise homeownership and reduce rental demand. Because the quantity of rental housing supplied is determined by the marginal rent earned on uncontrolled units, the overall effect of these two forces would raise average rents by 4.1\% but reduce the quantity of rental housing supplied by 0.7\%. Decontrolling also would reduce the rent of new tenants by 6.7\% compared to the benchmark. In figure \ref{fig:diagram_full_decontrol}, we provide a qualitative diagram that visually demonstrates some of the effects discussed in this and the previous paragraph. To save on space here, we provide detailed explanation in the notes at the bottom of the figure.

Overall, decontrolling would raise newborns' welfare by 0.26\%, but reduce average welfare by 0.96\%.\footnote{An increase in average welfare does not mean that all households gains; some lose, some gains, but the average increases. The reverse holds when average welfare decreases.} This is because newborns would benefit more from the lower initial rents under full decontrol than from the lower average rents in the benchmark. The opposite would be true for older households that have been under rent control for a while and have already began to benefit from it. 

\subsubsection{Distributional results} 

We have shown that full rent decontrol would increase the average welfare of newborns, but how would these gains be distributed? To answer this question, we break down newborns into seven percentiles according to their labor productivity and decompose the welfare gains and losses into three different effects: rent effect, house price effect, lump sum transfer effect. 

We present the distributional results in table \ref{tab:distributional_impact}. In the first three rows we present decomposed welfare changes. The bottom row of the table presents the total welfare changes--the sum of all the separate effects. $e_{1}$ represents the lowest labor productivity, and $e_{7}$ the highest. We begin by focusing on the first row in the table, which lists the welfare gains and losses arising only from changes in the rental market.\footnote{To obtain the rent effect, we hold house prices and lump sum transfers fixed at their levels in the benchmark long-run equilibrium. We then decontrol the rental market and set the free market rent to equal its value in the long-run equilibrium of the fully decontrolled rental market. We then run the model in partial equilibrium to obtain the direct effect of rents on welfare. To obtain the house price effect, we run another simulation, in which we repeat the exercise for the rent effect but also set the house price to its level in the long-run equilibrium of the fully decontrolled rental market. The lump sum transfer effect is the difference between the long-run equilibrium of the fully decontrolled market and the exercise in which we only change rents and house prices.} 

\paragraph{Rent effect} Newborn households with the lowest productivity gain substantially from rent decontrol. These households are born assetless renters and have the lowest incomes. They would have more after-tax income for non-housing consumption because of the lower rents for new tenants compared to the benchmark. These households would benefit from trading off lower rents early in life in exchange for higher average rents because their incomes are quite low early in life. The opposite would be true for the second and third lowest productivities, $e_{2}$ and $e_{3}$. As the second row in table \ref{tab:avg_years_in_rental} shows, these households expect to remain many years in the same rental unit, almost as many as households with $e_{1}$. However, their incomes are higher early in life. Thus, newborns with $e_{2}$ and $e_{3}$ value lower average rents more than lower rents upon moving into a new unit. With higher income, newborns with the top four productivities, $e_{4}-e_{7}$, would gain from rent decontrolling. As shown at the second row of table \ref{tab:avg_years_in_rental}, these households expect to remain significantly fewer years in the same rental units, and they are more likely to rent in the decontrolled market, because they expect to become homeowners rather quickly. This means that they do not benefit from rent control, and therefore would prefer to pay lower rents in the decontrolled market.

\paragraph{House price effect} In the second row of table \ref{tab:distributional_impact}, we present the welfare gains and losses that arise purely from the 0.8\% increase in house prices. All newborn households, regardless of productivity, would lose from the higher house price. This is because it would delay homeownership and reduce the house size they would expect to purchase. Since newborns with the lowest three productivies have higher expectations of remaining renters longer, they lose less than newborns with the highest four productivities, who expect to become homeowners more quickly and buy larger homes. 

\paragraph{Lump sum transfer effect} The third row of table \ref{tab:distributional_impact} lists the welfare gains and losses due to the lump sum transfer. As one would expect, the lump sum transfer would increase welfare across the board. However, newborns with lower productivities would gain more than newborns with higher productivities because the former have higher marginal utilities than the latter.

\paragraph{Total welfare changes} In the last row of table \ref{tab:distributional_impact}, we sum up all the welfare effects to arrive at the total welfare change. For newborns with $e_{1}-e_{3}$, the rent effect would outweigh the house price and lump sum transfer effects. For newborns with the top four productivities, all the effects would roughly cancel each other out, leading to negligible welfare gains or losses. 

Because the largest changes would occur in rental markets, and those changes would primarily affect low productivity newborns, the rent effect would have a first order effect on welfare. The lump sum transfer would have a second order effect on welfare because it would mostly benefit newborns with the lowest productivities who have the highest marginal utilities. House prices, in contrast, would mainly affect high productivity newborns who are more likely to become homeowners. Since these households have the lowest marginal utilities, the house price effect would have the least impact on welfare. 

To sum up, rent decontrol would reduce average housing affordability, but would also reduce rents of new tenants. Contrary to conventional wisdom, the latter effect would mainly benefit newborns whose incomes are lowest. Newborns with low- to-mid incomes would lose due to the rising average costs of renting, whereas high-income newborns would have negligible welfare gains and losses.

\subsection{Benchmark vs. full rent control}

In this experiment, we compare the long-run equilibrium under rent control without exemptions with the benchmark. This exercise highlights the impact of not allowing any exemptions on variables such as rental housing supplied, rents, and others aggregate outcomes. Because the distributional outcomes provide no new insights, we focus on aggregate outcomes here.  

The third column of table \ref{tab:aggregate_effects} presents changes in aggregate variables in the equilibrium under full rent control relative to the benchmark. With rental markets being fully controlled, renters who were renting in the free rental market in the benchmark would find it optimal to remain longer in the same rental units. This would also mean that will not as much as their incomes grow, and thus live in smaller units than in the benchmark. As a result, on average, renters would remain 0.85\% longer in the same units and demand 0.7\% smaller rental units compared to the benchmark. These two effects would reduce housing demand and thereby reduce equilibrium house prices by 0.3\%. The latter would mean that more renters would be able to purchase a house, and thus homeownership would increase by 0.3p.p.. 

With no exemptions, the marginal rent would become the average rent in the controlled market, which would be 6.77\% lower than the free market rent in the benchmark. This lower marginal rent would reduce the quantity of rental housing supplied by 1.4\%. Moreover, the decline in the marginal rent, would require higher initial rents for new tenants to clear the rental market, and this would push the average rent up by 0.4\%. Our results support the argument that rent control reduces the quantity of rental housing supplied and that allowing for exemptions negates some of the negative impact on the quantity supplied. Lastly, because full rent control would shut down the free rental market, where rents are lower than initial rents in the controlled market in the benchmark, it would force all new tenants to pay higher rents. This, together with higher average rents, would hurt virtually all households and reduce welfare by 0.50\%. Figure \ref{fig:diagram_full_control} presents a qualitative diagram to help the reader visualize a number of the mechanisms expounded on in this section. To reduce duplication, we provide further detailed explanations in the notes at the bottom of the figure.

At this point, a careful reader may begin to wonder how it is possible that homeownership would rise relative to the benchmark both when rental markets are fully decontrolled and when rental markets are fully controlled. The key to solving this puzzle is to notice that the increase in homeownership would arise due to different reasons. When markets are fully decontrolled, homeownership would rise despite house prices increasing, whereas when markets are fully controlled, homeownership would rise because house prices would decrease. Why is that? When markets are fully decontrolled, more renters would want to become homeowners because renting would become less desirable. This would increases the overall demand for housing and homeownership, and increase house prices. When markets are fully controlled, renters would be encouraged to move less, and thus demand smaller rental units. This would reduce overall housing demand and house prices. With lower house prices more renters would be able to attain the downpayment and ongoing mortgage payments, and thus homeownership would increase.

\subsection{Optimal degree of regulation}

Although we have shown that full decontrol is superior to the benchmark, a careful reader may begin to wonder whether some degree of rent control would be better than full decontrol. To answer this question, we perform a grid search to find the degree of rent control under which the ex-ante expected utility of a newborn household is maximized. We increase and decrease the share of the rental market that is mandated to be under rent control and compute the welfare associated with the long-run stationary equilibrium for each case. 

We find that no degree of rent control is optimal, which means that in the long-run welfare is highest under full decontrol. What explains this finding? When the degree of rent control declines, equilibrium rents in the controlled market must increase in order to reduce the mass of renters in the controlled market. This is an outcome of renters' sorting choice between the decontrolled market and the controlled market. Because renters with lower-incomes tend to remain in the same rental unit longer, they benefit more from living in a rent-controlled unit. Thus, in equilibria with a lesser degree of rent control, they sort into the controlled market, and end up being subjected to higher rents. Since these renters have the highest marginal utilities, their welfare loss also reduces overall welfare. 

\subsection{Transition from benchmark to full decontrol}

We have shown that full decontrol is superior in the long-run to all long-run equilibria in which part of the rental market is under control. But in many instances, policymakers care about the costs that would have to be incurred during the transition from one system to another. This raises the question: how long would it take for welfare gains to materialize? To answer this question we simulate the transition from the benchmark equilibrium to the long-run equilibrium under full decontrol. We assume that in period 1 households are completely surprised by the government permanently fully decontrolling rental markets, also known colloquially in the profession as a permanent "MIT shock". 

We present our results in the panels (a)-(c) of figure \ref{fig:transition_path}.\footnote{We do not show the homeownership rate as the changes in this variable are small and provide no additional insights.} In panel (a), the y-axis on the left side is the scaling for free market rents, where we normalize the benchmark to 1, and the y-axis on the right side is the scaling for welfare gains and losses. In panel (b), we show average rents, where normalize the benchmark to 1. In panel (c), the y-axis on the left side is the scaling for our house price, where we normalize the benchmark to 1, and the y-axis on the right side is the scaling for the lump sum transfers as a percentage of average labor income. In all panels, the x-axis represents periods since implementation.

As panel (a) shows, rents would decline 4.6\% upon implementation, which is 1.3p.p. more than their long-run equilibrium level. Since the homeownership rate remains roughly the same, this suggests that the decline in rents would be a consequence of renters reducing their units' sizes. Why would that happen? The key to understanding this, is to remember that the joint distribution of housing and financial assets across ages in the benchmark equilibrium is a function of households' choices under a high degree of rent control. Upon the removal of rent control, households endowments positions are fixed, and it would take them time to adjust their positions to reflect the optimal choices under full rent decontrol. Taking this into consideration, some renters would begin the period living in large units but would no longer benefit from rent control. These renters would see a sharp increase in their rents upon implementation, which would cause them to switch to smaller units, reducing the overall demand for rental housing, and thus decreasing equilibrium rents beyond their long-run equilibrium levels. This reduction in rents would be the primary force raising welfare. 

As the transition continues, however, households would anticipate the higher average rents and increase their savings in order to buy a house, which would eventually raises overall housing demand, the homeownership rate, and house prices to their higher long-run levels, as shown in panel (b). As house prices would rise, property tax revenues would too, and the government would be able to provide larger lump sum transfers as well.

Moreover, while waiting to purchase a home, renters would increase their average savings and thus enjoy greater capital income from their increased savings. With no additional monetary benefit to staying more periods in the same rental unit, this would encourage them to move more frequently to larger rental units as their incomes grow. This, in turn, would put upward pressure on the demand for rentals as well and raise rents over the next 30 periods slightly beyond their long-run levels, but eventually settling down to their long-run level. During the 30 periods during which rents would be above their long-run level, welfare gains would decline slightly below 0. However, after rents would settle at their long-run levels, welfare gains would increase to their long-run levels.

The key message from the transition exercise is that rents would decline rapidly upon implementation causing welfare gains to materialize immediately and remain positive almost all along the transition path.

\section{Conclusion}\label{sec:conclusion}

Many cities around the world have implemented or considered implementing rent controls that exempt newly-built units and allow landlords to reset rents upon tenant turnover \citep{fraser-1981}. These controls have been investigated empirically \citep{hayek-1929,fs-1946,fallis-1985,diamond-2019}, but little analysis has been done on the effects of rent control on equilibrium house prices, free market rents, and controlled rents, as well as its distributional impacts on household welfare \citep{fmn-2022}. 

We use a general equilibrium model with heterogeneous households to study tenancy deregulation, also referred to as rent control, in conjunction with exemptions, also referred to as the free market. We carefully calibrate our model to data on Toronto. We use various sources such as household-level Census data, unit-level Rental Market Survey data, and household-level Survey of Financial Security data. Using the calibrated model, we conduct four experiments. In the first experiment, we compare the long-run equilibrium under full rent decontrol to the benchmark. In the second, we examine the long-run equilibrium under full rent control relative to the benchmark. In the third, we search for the degree of rent control that produces the highest welfare in the long-run equilibrium. Lastly, we explore the transitional dynamics from the benchmark to the long-run equilibrium under full decontrol. In all of our experiments we keep the tax code fixed, and balance the budget using lump sum transfers. Using these experiments, we examine aggregate outcomes as well as distributional ones.

Our experiments reveal several key findings. First, fully decontrolling rental market would reduce rents paid by new tenants, but raise average rents. The former would benefit low-income, young households, while the latter would hurt renters that have lived in controlled rental housing for a while. Decontrolling would considerably reduce tenancy duration. The transition to full rent decontrol would exhibit an immediate decline of free market rents, which would raise welfare upon implementation, but also an increase in average rents. Second, full rent control would reduce the marginal rent received by the rental company, thus reducing the quantity of rental housing supplied and raising rents in the controlled market. As a result almost all households would lose. This finding suggests that exemptions mitigates the negative effect of full rent control on the rental housing supply. Third, there is no optimal degree of exemptions. In other words, we find that in the long-run full decontrol results in the equilibrium with the highest welfare gains. This finding is a result of renters' optimally sorting between free and controlled markets. When more rental units are exempted from rent control, equilibrium rents of controlled units would rise, leaving only those who benefit the most from rent control in that market. Since low-income households are least likely to become homeowners, they benefit the most from rent control. Thus, they would sort into the rent-controlled market, but face higher rents, resulting in welfare losses. 

Our study is among the first to model rent control with exemptions in a general equilibrium with heterogeneous households, which means that the general equilibrium literature on the topic has still room for improvement. We think that there are some important elements that should be included in future analyses once computational capacity is able to support it. Due to computational limitations, we abstracted from certain features such as endogenous labor supply and aggregate uncertainty, and this, of course limited the scope of our work. Abstracting from endogenous labor supply limited our ability to convincingly conduct policy experiments that consider whether rent subsidies paid for by higher labor income taxes could solve some of the negative impacts of rent control. Similarly, abstracting from aggregate uncertainty prevented us from exploring how the rigidity of rent control impacts household welfare during periods in which equilibrium rents in the free market exhibit substantial fluctuations around their long-run trends, but rents in the controlled market are more stable. It could be that this element is an important benefit of rent control that has yet to be fully understood.

%While the economic profession has a rather unanimous view about the harms of rent control, our analysis shows that rent control does have raison d'être, either for political reasons or fears of changes. We provide a general equilibrium analysis of current rent control with heterogeneous households. Our model shows that young and low-income households would be better off without rent control, and that rent control reduces rental housing supply and makes rental housing less affordable. It also shows that the transition to rent decontrol is not as scary as believed by immediately raising average rents but reducing rents of previously-uncontrolled units. 

%First, intermittent exemptions are commonly used as a policy tool, but equilibrium analyses have yet to capture this phenomenon properly to understand its implications. Our study partially advances the literature along this dimension by showing that permanent exemptions mitigate the negative effect of full rent control on the quantity of rental housing supplied, which causes rental prices to rise.

\bibliographystyle{aer}
\bibliography{reference_files/bib}

\begin{figure}[H]
  \centering
    \caption{Qualitative diagrammatic representation of rent path based on rental market choice and years to remain in the same rental}
    \vspace{-0.46cm}
    \label{fig:diagram_rent_path}
    \includegraphics[width=\textwidth]{figures/path_of_rent_diagram.pdf}
    \vspace{-1.4cm}
    \begin{fignote2}{1.0\textwidth}
      \textit{Notes:} The downward sloping red line shows rents paid by renters who live in controlled units as a function of their tenancy duration in the same units. The horizontal red line shows rents paid by renters who sort into the free market. In this market rent is independent of tenancy duration. Rents paid by renters renting in the free market are lower than rents paid by renters in the controlled market during their first years of, but higher than the average rent paid in the controlled market. Why is that? Consider the following thought experiment: if initial rents in the controlled market were equal or lower than free market rent, all renters would optimally sort into the controlled market. In order to shift some renters to the free market, renters in the controlled market must pay a premium initially. This means that renters anticipating a short stay in their rental unit would sort into the free market to avoid the premium. On the other hand, renters who expect to stay many years in the same rental would bite the bullet, pay the initial premium but enjoy lower average rents.      
    \end{fignote2}
\end{figure}

\begin{table}[H]
    \centering
    \caption{Assigned parameter values}
    \vspace{-0.5cm}
    \label{tab:calibration}
    \input{tables/calibration.tex}
    \vspace{-0.45cm}
    \begin{fignote2}{1.07\textwidth}
    \textit{Notes:} Table \ref{tab:calibration} we show the values of the calibrated parameters. Part (i) shows the external calibration for non-housing parameters. Part (ii) lists the external calibration of housing parameters. Part (iii) displays the external calibration of mortgage-related parameters. Part (iv) shows the internally assigned parameter values. The model matches all target almost exactly. Data for mortgages/housing value comes from the Survey of Financial Security. Data for total rent/total income, homeownership rate, average income tax rate, and renters with rent/income>30\% come from the 2016 Census. The rent-controlled market share is the average of the shares in the 1991-2016 Census waves. The elasticity of rental supply is estimated using data from the 2001-2006 Census waves on the quantity of bedrooms rented out and average real rents, as well as data from the MLS on average real growth in house prices. 
    \end{fignote2}
\end{table}

\begin{figure}[H]
  \centering
    \caption{Model performance on life-cycle statistics}
    \vspace{-0.5cm}
    \label{fig:life_cycle}
    \includegraphics[width=\textwidth]{figures/lifecycle_plots.pdf}
    \vspace{-0.45cm}
    \begin{fignote2}{0.8\textwidth}
      \textit{Notes:} In the two panels in figure \ref{fig:life_cycle}, we plot the performance of the model on life-cycle statistics relative to the data. In panel (a), we show the average homeownership rate of several age groups: 26-34, 35-44, 45-54, 55-64, and 65+. In panel (b), we also show the average of the LTV ratios for all homeowners in the same age group. The model output is in shown in blue and the data is in orange. Data for homeownership comes from the 2016 Census, and data for the LTV ratio comes from the 2016 Survey of Financial Security.       
    \end{fignote2}
\end{figure}

\begin{table}[H]
    \centering
    \caption{Aggregate non-targeted statistics}
    \vspace{-0.5cm}
    \label{tab:aggregate_non_targreted_statistics}  \input{tables/aggregate_non_targeted_statistics.tex}
    \vspace{-0.45cm}
    \begin{fignote2}{0.91\textwidth}
      \textit{Notes:} Table \ref{tab:aggregate_non_targreted_statistics} lists the performance of the model on aggregate non-targeted and distributional statistics relative to the data. The performance of the model is shown in the second column, while the third column has the estimates from the data. 
    \end{fignote2}
\end{table}

\begin{table}[H]
    \centering
    \caption{Tenure by labor productivity for all households in benchmark}
    \vspace{-0.5cm}
    \label{tab:tenure_by_labor_productivity}    \input{tables/tenure_by_productivity_benchmark.tex}
    \vspace{-0.45cm}
    \begin{fignote2}{0.96\textwidth}
      \textit{Notes:} Table \ref{tab:tenure_by_labor_productivity} presents tenure by labor  productivity for all renters. $e_{1}$ represents the lowest labor productivity, and $e_{7}$ the highest. In the brackets next to each productivity, we denote the percentage of households that have that productivity. For example, 1.72\% of households have productivity $e_{1}$. The columns under each labor productivity show the distribution of households with that productivity across tenure. For example, the first row of the $e_{1}$ column says that 76.14\% of households with productivity $e_{1}$ are renters in the controlled market. The second row says that 2.01\% of households with $e_{1}$ are renters in the free market. The third row says 8.5\% of households with $e_{1}$ are owners with a mortgage, and the last row says that 13.35\% are homeowners without a mortgage.
    \end{fignote2}
\end{table}

\begin{table}[H]
    \centering
    \caption{Aggregate effects}
    \vspace{-0.5cm}
    \label{tab:aggregate_effects}
    \input{tables/aggregate_effects.tex}
    \vspace{-0.45cm}
    \begin{fignote2}{0.66\textwidth}
      \textit{Notes:} The second column of table \ref{tab:aggregate_effects} compares the benchmark equilibrium to the long-run equilibrium under full decontrol. The third column of table \ref{tab:aggregate_effects} compares the benchmark equilibrium to the long-run equilibrium under full control. "--" means that the statistic cannot be computed.
    \end{fignote2}
\end{table}

\begin{figure}[H]
  \centering
    \caption{Qualitative diagrammatic representation of the impact of full rent decontrol}
    \vspace{-0.37cm}
    \label{fig:diagram_full_decontrol}
    \includegraphics[width=\textwidth]{figures/full_rent_decontrol_diagram.pdf}
    \vspace{-0.8cm}
    \begin{fignote2}{1.0\textwidth}
      \textit{Notes:} In figure \label{fig:diagram_full_decontrol} we provide a qualitative diagrammatic representation of the impact of full rent decontrol relative to the benchmark. Full rent decontrol would shift both the rental demand curve and the rental supply curve to the left, from $\mbox{Rental demand}_{1}$ to $\mbox{Rental demand}_{2}$, and from $\mbox{Rental supply}_{1}$ to $\mbox{Rental supply}_{2}$, respectively. Two opposing forces would affect the rental demand curve. On the one hand, full rent decontrol would reduce the incentive to rent, increase homeownership, and thus reduce rental demand. On the other hand, full rent decontrol would increase savings to purchase a house, increase capital income, and strengthen the incentive of the remaining renters to move to larger rental as their incomes grow. The former effect would dominate the latter, causing demand to shift to the left. Similarly to demand, supply would also shift to the left as a result of an increase in equilibrium house prices. This shift can be seen through equation (\ref{eq:rental_price}). The benchmark equilibrium shown has quantity, $X^{1}$, free market rent, $p^{1}_{1}$, average controlled rent $\overline{p}_{0}$, and average rent $\overline{p}$. The equilibrium under full rent decontrol has quantity $X^{2}<X^{1}$, free market rent $p^{2}_{1}$, and average rent, which is the same as the free market rent in the new equilibrium, $p^{2}_{1}>\overline{p}$. Why would $p^{2}_{1}>\overline{p}$, but $X^{2}<X^{1}$? The average rent would increase relative to the benchmark because the increase in rents charged on previously-controlled units in the benchmark, $\underline{X}_{0}$, would dominate the reduction rents obtained in the free market when the share of the free rental market in the benchmark is small. The quantity of rental housing supplied would decline relative to the benchmark because free market rents would decline. This would lead to a lower marginal rent commanded by the rental company, and thus reduce its quantity supplied. 
    \end{fignote2}
\end{figure}

\begin{table}[H]
    \centering
    \caption{Avg. years in same rental unit by labor productivity for all renters}
    \vspace{-0.5cm}
    \label{tab:avg_years_in_rental}
    \input{tables/years_in_same_rental.tex}
    \vspace{-0.45cm}
    \begin{fignote2}{1.04\textwidth}
      \textit{Notes:} Table \ref{tab:avg_years_in_rental} reports tenure length of renters in the model broken down by labor productivity. $e_{1}$ represents the lowest labor productivity, and $e_{7}$ the highest. In the brackets next to each productivity, we denote the percentage of households that have that productivity. For example, 1.72\% of households have productivity $e_{1}$. The first row shows the duration a renter stays in the same unit in the benchmark where rental markets are partially regulated. The second row shows the duration when we simulate the model without any regulation.
    \end{fignote2}
\end{table}

\begin{table}[H]
    \centering
    \caption{Newborns' distributional impact of full decontrol by labor productivity}
    \vspace{-0.5cm}
    \label{tab:distributional_impact}
    \input{tables/distributional_impacts.tex}
    \vspace{-0.45cm}
    \begin{fignote2}{0.95\textwidth}
      \textit{Notes:} Table \ref{tab:distributional_impact} displays welfare effects for newborn households, comparing full rent decontrol to the benchmark. $e_{1}$ represents the lowest labor productivity, and $e_{7}$ the highest. In the brackets next to each productivity, we denote the percentage of newborn households that have that productivity in their age. For example, 1.25\% of newborn households have productivity $e_{1}$. Note that the percentage of newborns in each productivity group is different than the percentage of all households in each productivity. This is because the intra- and inter-generational labor productivity processes imply that the stationary ergodic distributions of labor productivities across all ages is different from that distribution of productivities of the youngest cohort. The first row isolates the welfare impact due to the change in rents. The second row isolates the impact due to the change in house prices. The third row provides the welfare effect due to the lump sum transfer. For each column, representing a labor productivity group, summing the numbers in bottom three rows, yields the number in the last row, which is the total welfare gain.
    \end{fignote2}
\end{table}

\begin{figure}[H]
  \centering
    \caption{Qualitative diagrammatic representation of the impact of full rent control}
    \vspace{-0.36cm}
    \label{fig:diagram_full_control}
    \includegraphics[width=\textwidth]{figures/full_rent_control_diagram.pdf}
    \vspace{-0.45cm}
    \begin{fignote2}{1.0\textwidth}
      \textit{Notes:} In this figure we present a qualitative diagram representing the impact of full rent control relative to the benchmark. Full rent control would shift the rental demand curve to the left, from $\mbox{Rental demand}_{1}$ to $\mbox{Rental demand}_{2}$, whereas the rental supply curve would shift to the right, from $\mbox{Rental supply}_{1}$ to $\mbox{Rental supply}_{2}$. As in the case of full rent decontrol, there would be two opposing forces affecting the shift in rental demand. Full rent control would increase the incentive to stay longer in the same rental unit, and thus decrease the average size of a rental unit demanded. This in turn would reduce overall housing demand, and cause equilibrium house prices to slightly decline. With lower house prices, more renters would be able to purchase a house, and thus homeownership would rise, reducing the demand for rental housing. Quantitatively, the former effect would dominate the latter, causing demand to shift to the left. In contrast to demand, supply would also shift to the right as a result of an decrease in equilibrium house prices. The direction of the supply shift can be gleaned from equation (\ref{eq:rental_price}). The benchmark equilibrium shown has quantity, $X^{1}$, free market rent, $p^{1}_{1}$, average controlled rent $\overline{p}^{1}_{0}$, and average rent $\overline{p}$. The equilibrium under full rent control has quantity $X^{2}$, average controlled rent $\overline{p}^{2}_{0}$, and average rent, which is the same as the average controlled rent in the new equilibrium, $\overline{p}^{2}_{0}>\overline{p}$. In the benchmark, the quantity supplied is determined by the free market rent, $p^{1}_{1}$ which is the marginal rent. In contrast, under full rent control, the quantity supplied is determined by the average controlled rent, $\overline{p}^{2}_{1}$. Since $\overline{p}^{2}_{1}<p^{1}_{1}$, the quantity supplied under full rent control, $X^{2}$, is lower than the quantity supplied in the benchmark, $X^{1}$.
    \end{fignote2}
\end{figure}

\begin{figure}[H]
  \centering
    \caption{Transition path}
    \vspace{-0.5cm}
    \label{fig:transition_path}
    \includegraphics[width=\textwidth]{figures/transition_plots.pdf}
    \vspace{-0.45cm}
    \begin{fignote2}{0.99\textwidth}
      \textit{Notes:} Figure \ref{fig:transition_path} plots the transition path of the model from the benchmark to the long-run equilibrium under full rent decontrol. In panel (a), the y-axis on the left side is the scaling for free market rents, where we normalize the benchmark free market rent to 1. The y-axis on the right side is the scaling for welfare gains and losses. Because welfare gains are 0 in period 0, which is the period before full decontrol, we plot welfare gains for period 1 and onward. In panel (b), we plot the average rent in the counterfactual relative to the benchmark, where we normalize the benchmark price to 1. In panel (c), the y-axis on the left side is the scaling house prices, where we normalize the benchmark price to 1. The y-axis on the right side is the scaling for the lump sum transfers as a percentage of average labor income. In all panels the x-axis corresponds to periods.
    \end{fignote2}
\end{figure}

\end{document}

%The increasing cost of housing in major cities is one of the greatest economic challenges of our time. In 2016, the share of renters paying over 30\% of their income on housing has risen by 25\% since 1990 to nearly half of all renters \citep{JCHS-2018}. In Canada, the share of renters spending over 30\% of their income on rent has ?? between ?? and 2016. The problem is even worse in major cities. In Toronto and Vancouver.... In response to these trends, rent control has recently returned to the limelight as a policy tool to increase housing affordability (need references here). 

%The second moment is the ratio of aggregate mortgage debt to aggregate housing value of homeowners, which we estimate from the 2016 wave of the Survey of Financial Security (SFS). We calibrate the discount factor to match this moment. The discount factor is important because it determines how much renters are willing to trade off lower rents early in life at the expense of higher average rents over their life-cycle. As we show shortly below, the main trade off between universal rent control and a universally deregulated rental market is that in the former rents are higher early in life but lower on average compared to the latter. 

%Before retirement, owners have access to mortgage debt, $m$, which is limited to $m \leq \lambda_{m}(ph) p h$, where $\lambda_{m}(ph)<1$, and depends on the value of the house. We set the interest rate on mortgage debt to $r_{m}=r+\Delta r$, where $r$ is the return on financial wealth, and $\Delta r$ is the interest premium paid on mortgage debt \citep{sommeretal-2013,sommersullivan-2017,kaplanetal-2019}. Following \citet{gervais}, mortgage debt is fully repaid at the beginning of each period. Because in Toronto households cannot deduct their mortgage interest or property tax payments from their taxable income, they cannot do so in the model either.

%Because in our calibrated model households' discount the future more than the present, $\beta<1$, in equilibrium, free market rent must always be strictly higher than the average controlled market rent. To see why, suppose that the free market rent was equal to the average rent in the controlled market. Then, households expecting to pay the same rent on average in both markets would sort into the free market because they would pay initially in the free rental market. This would raise the free market rent and lower the average rent in the controlled market.

%As well, the model's percentage of owners with a housing costs to income ratio above 30\% and 50\% is also similar to the same statistics we estimate using 2016 SFS data.

%This is due to the fact that many households in Toronto have high LTV ratios at origination, but the high real growth rate in house prices quickly reduces their mortgage debt relative to the value of their homes, showing up as lower LTV than in our model. In fact, a household purchasing a home with a LTV ratio of 80\% at origination and does not pay any of the principal in the first 10 years, would have a LTV ratio of only 51\% at the end of the 10 years as a result of the 4.46\% real growth in house prices in Toronto estimated in the data. 

\begin{comment}
We use 2016 census data of Toronto CMA to compute several key parameters. Homeownership rate is 67\%. Total rent to total income and renters with rent/income higher than 30\% (calculate in VDL). In 2016, the MLS average price is \$767,336 and average rent of primary apartments is \$1,240 from Rental Market Survey. The average rent of secondary apartments is \$1,901 and the average rent of other dwelling types other than apartments is \$1,583. Taking the average rent of primary and secondary markets, the ratio of rent over house price is 2.46\%. The HPI for apartments in 2016 is \$359,618 and average rent of apartments is \$1570.5. The ratio of rent over apartment price is 5.24\%. The HPI for apartments in 2006 is \$329,650 and average rent of apartments is \$1021. The ratio of rent over apartment price is 3.72\%.
From CMHC (2018), the elasticity between housing starts and house prices in Toronto from 1992 to 2016 is 0.35. We also collected property tax rates from different municipalities in Toronto CMA and computed the average property tax rate, 0.67\%. The average real growth rate of house prices in Toronto from 2006 to 2016 is 4.46\% using MLS average house prices (that's more workable). The average growth rate of house prices in Toronto from 2006 to 2021 is 4.04\%. I'll take the 4.46\% and see what happens!. The ratio of total rent to total income before tax for renters in Toronto CMA in 2016 is 24.29\%. The share of renters with rent/income over 30\% is 45.16\% in Toronto in 2016. The average household employment income in Toronto in 2016 is \$74,722. 
\end{comment}

%Next steps:
\begin{enumerate}
    \begin{comment}
    \item Average income tax rate for Toronto - aggregate taxes paid over aggregate gross income. 
    I used aggregate numbers. The income tax rate is 19.46\%. Does that include CPP and EI contributions? CPP and EI are not included. (I see EI contributions are 1.88\% and CPP are 4.95\% in 2016). So about 26\% if those are not included. Still really low...I think we are missing something. I think you should limit it to households aged above 22 and below 65. We should re-estimate $\psi$ restricting the ages to 22-65 as well. Even for the age from 22 to 65, the income tax rate is 19.99\%. No big difference. $\psi$ is 0.9278 and constant 0.6718.
    \item From SFS (2016) LTV 26-34 - 44.1\%, 35-44 - 40.1\%, 45-54 - 25.7\%, 55-64 - 16.9\%, 65+ 4.8\%
    \item From SFS (2016) top 1\% - \%, top 10\% - \%, top 20\% - \%, bottom 50\% - \%, Gini -
    \item Maybe someone did similar work for intra-generational log(earnings), https://www.bankofcanada.ca/wp-content/uploads/2020/11/swp2020-46.pdf? https://www23.statcan.gc.ca/imdb/p2SV.pl?Function=getSurvey&Id=31183 - yes, Ihttps://www.sciencedirect.com/science/article/abs/pii/S1094202512000397
    \item SR: Increase in real income of households between ages 25-55 (and 65).
    The growth rate for 25-65 is 104\%
    The growth rate for 25-55 is 147.86\%
    average income for 24-26 year-old households: \$56,151.55
    average income for 54-56 year-old households: \$137,802, and 64-66 year-old households:\$113,520.8.
    Interesting. So declining after 55....I guess we should have an up and then down. I made the function go up until age 55 and then go down until 65.
    \item Fraction of Toronto's rental market not under regulation.
    From RMS, the fraction of uncontrolled units in Toronto in 2016 is 2.62\%. It is the fraction of rental units built after 1992. Interesting - very low! It is basically full regulation.
    \item SR: Lin, what source did you use to get this from "The HPI for apartments in 2016 is \$359,618 and average rent of apartments is \$1570.5. The ratio of rent over apartment price is 5.24\%"?
    HPI comes from CREA.
    \item Homeownership rate for each age 25-86 from the Census.
    In 2016, homeownership rates for 25-34 years 42\%, for 35-44 years 64.3\%, for 45-54 years 72.2\%, for 55-64 75\%, and for 65 years and over 75.1\%.
    \item SR: I will compute the LTV for each age from 25-86.
    \item SR: I will compute wealth inequality from SFS.
    \item SR: https://www.oecd.org/els/public-pensions/PAG2021-country-profile-Canada.pdf
    \item SR: Lin, you'll see in the text there are some places where I put ?? below. Can you add some description of how and where you got the statistic from, for all these ?? 
    \item SR: we need the fraction of gross income from all sources spent on rent (net of any rent subsidy) at the bottom 2\%, 3\%,....,10\% of the income distribution. Might not need this because I can calibrate the smallest rental so that no one spends over 75\% of their gross income on rent, which I think may be a reasonable thing to do. Otherwise, our welfare results may be too dependent on this number.
    \item SR: mortgage interest rate premium relative to risk free rate? Currently using 1\% based on US data, but would be better if we had this data for Canada.
    \item SR: we need an estimate of the elasticity of rental housing supply to calibrate $\theta$. LZ: model to estimate elasticity
    $log(starts)=\alpha log(rent)+controls+\epsilon$
    \item SR: for completeness in the non-targeted stats we should also show how many homeowners spend over 30\% and 50\% of their income on housing (property taxes+interest+principal), excluding those spending over 75\%.
    \item SR: is it possible to compute the \% of renters spending over 30\%, 50\% of their income on housing? And in these calculation exclude households spending over 75\% of their income on rent (I think those are quite odd cases).
    \item I think that in order to motivate ourselves in the beginning, it would be good to show statistics of over time (from 1976-2016 - or whatever far back we can go) on housing affordability in Canada as a whole, and in major centre like Toronto, Vancouver and Montreal.
    \end{comment}
    
\end{enumerate}